% ============================================================
%  Epistemic Reduction and Admissible Projection
%  Anderson's Epistemology as a Theorem of RSVP Field Theory
%
%  Author : Flyxion
%  Engine : LuaLaTeX
%  Font   : TeX Gyre Pagella
% ============================================================

\documentclass[12pt, a4paper, oneside]{book}

% ── Engines & Fonts ─────────────────────────────────────────
\usepackage{fontspec}
\setmainfont{TeX Gyre Pagella}
\usepackage{microtype}

% ── Mathematics ─────────────────────────────────────────────
\usepackage{amsmath, amsthm, amssymb, mathtools}
\usepackage{thmtools}

% ── Page Geometry ───────────────────────────────────────────
\usepackage[
  top=30mm, bottom=30mm,
  left=32mm, right=32mm,
  headheight=15pt
]{geometry}

% ── Headers & Footers ───────────────────────────────────────
\usepackage{fancyhdr}
\pagestyle{fancy}
\fancyhf{}
\fancyhead[L]{\small\itshape\leftmark}
\fancyhead[R]{\small\thepage}
\renewcommand{\headrulewidth}{0.3pt}

% ── Section Styling ─────────────────────────────────────────
\usepackage{titlesec}
\titleformat{\chapter}[display]
  {\normalfont\Large\scshape}
  {\chaptertitlename\ \thechapter}{1em}
  {\Huge\bfseries}
\titlespacing*{\chapter}{0pt}{40pt}{30pt}

% ── Hyperlinks ──────────────────────────────────────────────
\usepackage[
  colorlinks=true,
  linkcolor=black,
  citecolor=black,
  urlcolor=black
]{hyperref}

% ── Theorem Environments ────────────────────────────────────
\declaretheoremstyle[
  headfont=\normalfont\bfseries,
  bodyfont=\normalfont\itshape,
  spaceabove=10pt, spacebelow=6pt
]{thmstyle}

\declaretheoremstyle[
  headfont=\normalfont\bfseries,
  bodyfont=\normalfont,
  spaceabove=10pt, spacebelow=6pt
]{defstyle}

\declaretheoremstyle[
  headfont=\normalfont\itshape,
  bodyfont=\normalfont,
  spaceabove=6pt, spacebelow=4pt
]{rmkstyle}

\declaretheoremstyle[
  headfont=\normalfont\bfseries,
  bodyfont=\normalfont\itshape,
  spaceabove=10pt, spacebelow=6pt,
  headpunct={.}
]{conjstyle}

\declaretheorem[style=thmstyle,   numberwithin=chapter]{theorem}
\declaretheorem[style=thmstyle,   sibling=theorem]{proposition}
\declaretheorem[style=thmstyle,   sibling=theorem]{lemma}
\declaretheorem[style=thmstyle,   sibling=theorem]{corollary}
\declaretheorem[style=defstyle,   sibling=theorem]{definition}
\declaretheorem[style=defstyle,   sibling=theorem]{example}
\declaretheorem[style=rmkstyle,   sibling=theorem]{remark}
\declaretheorem[style=conjstyle,  sibling=theorem]{conjecture}

% ── Custom Operators & Macros ───────────────────────────────
\newcommand{\Adm}{\mathcal{A}}
\newcommand{\X}{\mathcal{X}}
\newcommand{\M}{\mathcal{M}}
\newcommand{\C}{\mathcal{C}}
\newcommand{\proj}{\pi}
\newcommand{\projA}{\pi_{\!\Adm}}
\newcommand{\dH}{d_H}
\newcommand{\SA}{{\sim_{\!\Adm}}}
\newcommand{\Nov}{\mathrm{Nov}}
\newcommand{\CC}{\mathrm{CC}}
\newcommand{\Ent}{S}
\newcommand{\phiCe}{\phi_{\C}^{\varepsilon}}
\newcommand{\vel}{\mathbf{v}}
\DeclareMathOperator{\argmin}{arg\,min}
\DeclareMathOperator{\colim}{colim}

% ── Miscellaneous ───────────────────────────────────────────
\usepackage{csquotes}
\usepackage{booktabs}
\usepackage{graphicx}
\usepackage{xcolor}
\usepackage{epigraph}
\setlength{\epigraphwidth}{0.60\textwidth}
\renewcommand{\epigraphflush}{flushleft}

% ============================================================
\begin{document}

\begin{titlepage}
  \begin{center}
    \vspace*{36mm}
    {\Huge\bfseries\scshape Epistemic Reduction\\[8pt]
    and Admissible Projection}\\[14pt]
    {\Large Anderson's Epistemology as a Theorem\\[4pt]
    of RSVP Field Theory}\\[44mm]
    {\large Flyxion}\\[8pt]
    {\normalsize Independent Researcher}\\[64mm]
    {\small Standard Galactic Edition}
  \end{center}
\end{titlepage}

\thispagestyle{empty}
\vspace*{\fill}
\noindent
Copyright \textcopyright\ Flyxion. All rights reserved.
No part of this monograph may be reproduced without attribution.\\[6pt]
\noindent
Monica Anderson's original texts are cited and quoted with full attribution.
The formal developments, theorems, and RSVP machinery are original
contributions of the present author.\\[6pt]
\noindent
Typeset with Lua\LaTeX\ using TeX Gyre Pagella.
\clearpage

\tableofcontents
\clearpage

% ============================================================
\chapter*{Preface}
\addcontentsline{toc}{chapter}{Preface}
\markboth{Preface}{Preface}

\epigraph{%
  You cannot Reason about that which you do not Understand.%
}{\textit{Monica Anderson}}

This monograph has a precise thesis: Monica Anderson's epistemological
program is derivable from the geometry of admissible projection and
RSVP field theory, not merely compatible with it.

The derivation proceeds through the admissibility quotient construction.
Once a trajectory space $\X$ and a continuous admissibility operator
$\Adm$ are in hand, Anderson's distinction between Understanding and
Reasoning emerges automatically as the difference between constructing
and operating within a representational manifold. Her Priority Thesis is
a structural tautology about the dependency between projections and
operators-on-projections. Her Salience criterion becomes admissibility
sensitivity. Her Corpus Congruence becomes the level-set structure of
a smoothed familiarity potential field. Her critique of classical AI
becomes the Projection Supremacy Failure Theorem. Her Wisdom Salon
becomes a sheaf-theoretic colimit whose global section is guaranteed
by the descent axiom. Calvin's uphill river becomes a monotone chain
of admissibility quotients with decreasing mutual information.

None of these require Anderson's framework as an assumption.

The monograph is mathematically candid about the status of each result.
Six categories of claim appear. Theorems are proved from the stated
hypotheses. Propositions are theorems with lighter hypotheses or shorter
proofs. Lemmas are auxiliary results used in other proofs. Corollaries
follow immediately from preceding results. Conjectures are plausible
claims for which a proof route is described but not completed; they are
explicitly labelled and the missing ingredients are identified. Remarks
interpret and contextualize.

The principal remaining conjectural element is the bounded curvature
characterization of admissibility-compatible learning: the Local
Extension Stability Theorem provides a genuine one-directional result,
but the biconditional that the previous draft asserted has been
downgraded to a conjecture pending further geometric development.
Readers who find the conjecture interesting will find the missing
ingredients identified clearly in Chapter~\ref{ch:learning}.

Anderson is named in the subtitle as an intellectual predecessor. This
monograph provides the proofs she did not supply, with the honesty
their subject requires.

\clearpage

% ============================================================
\part{Anderson's Epistemological System}

\chapter{Intelligence, Understanding, and Reasoning}
\label{ch:intelligence}

\epigraph{%
  Intelligence = Understanding + Reasoning.%
}{\textit{Monica Anderson, \emph{Why AI Works}}}

\section{The Central Decomposition}

Anderson's framework opens with a functional decomposition of
intelligence into two irreducible faculties. Understanding is
subconscious, intuitive, model-free, holistic, expensive, and fallible
in a characteristic way: its errors are graceful near-misses rather than
catastrophic failures at arbitrary inputs. Reasoning is conscious,
logical, model-based, cheap per inference step, and conditionally
infallible given correct input within a specified model space.

The central structural claim, stated explicitly by Anderson and repeated
in multiple forms, is that Understanding is \emph{prior to} Reasoning.
One cannot reason about a domain one has not yet understood. The converse
does not hold. This asymmetry Anderson calls the Priority Thesis.
Part~III derives it as Theorem~\ref{thm:priority}.

\section{Properties of the Two Faculties}

Understanding is subconscious (no introspective access to its operation),
intuitive (rapid and parallel, answering without deliberation), model-free
(no a priori symbolic representation of the problem domain), holistic
(context is taken as a whole), expensive (the majority of a brain's
processing budget), and fallible. Reasoning is conscious, sequential,
model-based, cheap per inference step, and conditionally infallible.
The guarantee is real but narrow: it applies only inside the model space
from which context has been deliberately stripped.

Anderson draws on Kahneman's System~1 / System~2 distinction but diverges
in one important respect. Kahneman treats the two systems as interacting
processes within a single cognitive architecture. Anderson treats them as
requiring fundamentally different computational paradigms with different
failure modes.

\section{The Priority Thesis}

\begin{proposition}[Anderson's Priority Thesis, informal]
\label{prop:priority-informal}
Understanding is strictly prior to Reasoning. Understanding does not
presuppose Reasoning. Reasoning presupposes Understanding.
\end{proposition}

This generates three consequences for AI methodology. First, any
architecture that attempts to produce understanding through symbolic
reasoning inverts the dependency. Second, twentieth-century AI (build a
model, reason over it) attacked the second step while bypassing the first.
Third, the correct question about AI architecture is epistemological.

Part~III proves Proposition~\ref{prop:priority-informal} as Theorem~\ref{thm:priority}.

% ──────────────────────────────────────────────────────────────
\chapter{Models, Reductionism, and the Pathology of Projection}
\label{ch:models}

\epigraph{%
  Model Based Problem Solving is the greatest invention our species
  has ever made.%
}{\textit{Monica Anderson, \emph{Our Greatest Invention}}}

\section{What a Model Is}

\begin{definition}[Model, after Anderson]
A Model is a simplification of reality produced by discarding context
in order to make a fragment of the world tractable to formal
manipulation. Models include scientific theories, equations, statistical
models, computer programs, and naive personal simplifications.
\end{definition}

\begin{definition}[Reductionism, after Anderson]
Reductionism is the use of Models.
\end{definition}

Models are the most productive intellectual tools humanity has produced.
A Model stripped of all domain-specific detail can be applied
by anyone trained in it, independently of whether they have any direct
experience with the original problem context.
Newton's $F = MA$ works in every context in which mass, force, and
acceleration are well-defined precisely because every other detail has been
discarded. The scientific enterprise of the past four centuries is largely
a library of such reusable compressions.

To create a Model is to perform a compression. The scientist surveys a
domain, identifies the features that determine the phenomenon of interest,
discards the rest, and publishes the result as a context-free representation
that others can apply without repeating the original investigation.
Anderson's central observation is that this compression is not a
description of a process internal to science.
It is a pre-scientific operation that science presupposes.
The decision about which features to retain and which to discard must
be made by an agent who already understands the domain at a level that
precedes the Model.
The Model is the output of Understanding, not its vehicle.

\section{The Model--Projection Correspondence}

The relationship between Anderson's Models and the formal machinery of
Part~II is exact rather than metaphorical.

\begin{proposition}[Models are Projections]
\label{prop:models-are-projections}
In the framework of Part~II, a Model in Anderson's sense corresponds to a
projection $\proj : \X \to \M$ from a trajectory space $\X$ to a
representational manifold $\M$.
The act of Model creation is the act of specifying $\proj$.
The context-free content of the Model is the structure of $\M$.
The discarded context is the information in the fibers $\proj^{-1}(m)$
for each $m \in \M$.
\end{proposition}

This correspondence makes Anderson's claims about Models testable as
claims about projections. Model reusability becomes the claim that $\M$
supports correct reasoning independently of which $x \in \proj^{-1}(m)$
was the original input. Model applicability becomes the claim that
$\proj^{-1}(\M)$ covers the region of $\X$ where the Model is being used.
Model failure becomes the claim that the input $x$ lies in
$\X \setminus \proj^{-1}(\M)$.

\section{Compression, Utility, and the Utility Threshold}

The compression effected by a Model is not a loss but a gain, under the
right conditions. A Model that compresses a $10^6$-dimensional state
space to a three-parameter equation does not impoverish the analyst;
it enables the analyst to solve problems that would be computationally
impossible without the compression. The utility of a Model is precisely
a function of its compression ratio together with the accuracy of the
admissibility structure it preserves.

\begin{definition}[Model Utility]
\label{def:model-utility}
The utility of a projection $\proj : \X \to \M$ for a class of tasks
$\mathcal{T}$ is the product of its compression (the reduction in
representational complexity $C(\X) - C(\M)$) and its admissibility
fidelity (the degree to which $\Adm_\M(\proj(x))$ approximates
$\proj(\Adm(x))$ for $x$ drawn from the task distribution).
A Model is useful when high compression and high admissibility fidelity
can be achieved simultaneously.
\end{definition}

The scientific enterprise can be understood as the sustained search for
projections with high utility: cases where the world happens to admit
compressions that preserve enough admissibility structure to support
reliable inference and action. Physics is the discipline where such
compressions are most available. The life sciences, social sciences, and
natural language are domains where they are most difficult to find,
which is why those fields resist the kind of mathematical formalization
that has been so productive in physics.

\section{The Pathology: Projection Supremacy}

The pathology Anderson calls Reductionism is not the use of Models.
It is a specific failure mode that occurs when Models are mistaken
for the reality they compress.
Anderson's target is not compression as such but the cognitive error of
treating the projection as complete: acting as though $\M = \X$.

This error has a precise name in the framework of Part~II:
Projection Supremacy (Definition~\ref{def:projsupremacy}).
A system in a state of Projection Supremacy applies its reasoning
operator $R : \M \to \M$ to inputs $x$ with non-zero projection error
$\dH(\Adm(x), \Adm_\M(\proj(x))) > 0$, without any mechanism for
detecting or compensating for this error.

The error is not that the system has a Model.
The error is that the system has no representation of the gap between
the Model and the world.

\begin{proposition}[The Model--Projection Supremacy Chain]
\label{prop:mps-chain}
The following sequence describes the generation and eventual failure of
a Reductionist cognitive system.
\begin{enumerate}
  \item \textbf{Model creation.} A scientist performs Epistemic Reduction:
    $\proj : \X \to \M$ is constructed, discarding non-salient features.
    The result is a useful compression with well-defined applicability
    conditions.
  \item \textbf{Institutionalization.} The Model is published,
    taught, and applied. Users learn $R : \M \to \M$ (the reasoning
    procedures of the field) without learning the construction of
    $\proj$. The original Understanding that produced the Model
    is separated from the Model itself.
  \item \textbf{Scope inflation.} The Model is applied in contexts
    increasingly distant from its original derivation. The applicability
    conditions are either unknown to the users or disregarded.
    The projection error $\dH(\Adm(x), \Adm_\M(\proj(x)))$ grows.
  \item \textbf{Projection Supremacy.} The system has no representation
    of the projection error. It applies $R$ uniformly, producing
    confident outputs even in regions where $\proj^{-1}(\M)$ does not
    cover $x$.
  \item \textbf{Brittleness.} Catastrophic failure occurs at the edges
    of competence by Theorem~\ref{thm:brittleness}. The failure is
    unmodulated because the system has no internal signal for the
    growing discrepancy between $\M$ and $\X$.
\end{enumerate}
\end{proposition}

\begin{remark}
Proposition~\ref{prop:mps-chain} is not a theorem; it is a phenomenological
description of a pattern that Proposition~\ref{prop:barrier} and
Theorem~\ref{thm:brittleness} together explain.
The sequence makes the transition from Anderson's critique of
classical AI to the RSVP formal treatment legible: Anderson is describing
the middle three steps of the chain (institutionalization, scope inflation,
Projection Supremacy) in operational terms.
The RSVP framework provides the geometry that explains why the fifth
step is structurally inevitable once the fourth is in place.
\end{remark}

\section{Why Models Cannot Produce Understanding}

Anderson's argument that Models cannot produce Understanding is the
circularity corollary stated below.
But in the expanded framework above, a more precise version can be given.

\begin{proposition}[Irreducibility Barrier]
\label{prop:barrier}
There exists a class of problem domains, which Anderson calls
Bizarre Domains, in which no projection $\proj : \X \to \M$ can be
simultaneously tractable ($\dim \M \ll \dim \X$) and admissibility-faithful
(projection error below any fixed threshold), because the admissible
future structure of $\X$ is not well-approximated by any low-dimensional
quotient. Reductionist methods are structurally blocked in these domains.
\end{proposition}

\begin{proof}[Proof sketch]
A domain is Bizarre in Anderson's sense when the phenomena of interest
are emergent: they depend on the entire configuration of the system and
not on any low-dimensional projection of it. Formally, this means the
equivalence classes of $\SA$ are singletons or near-singletons:
almost every state has a unique admissible future set. In this case the
admissibility quotient $\M_0 = \X/\SA$ is essentially isomorphic to $\X$
itself, and no compression of substantial ratio is admissibility-faithful.
\end{proof}

\begin{corollary}[Model-Based AI Circularity]
\label{cor:circularity}
Any AI architecture that attempts to produce Understanding through
hand-crafted Models is self-defeating. Understanding is the construction
of $\proj : \X \to \M$. A hand-crafted Model is a fixed $\proj$,
specified in advance of the data. By Theorem~\ref{thm:feyerabend},
any feature discarded by this fixed $\proj$ cannot become salient
through reasoning operations internal to $\M$.
The architecture can therefore never discover admissibility-relevant
features that its designers did not already know mattered.
\end{corollary}

\begin{remark}
The corollary is stronger than a circularity argument.
It is an impossibility result: a hand-crafted Model not only fails to
generate Understanding for the features it retains; it structurally
prevents the discovery of salience for the features it discards.
This is the information-theoretic content of Anderson's claim that
``models cannot create understanding'': not only is there a circular
dependency, but the projection forecloses the discovery of its own
inadequacy.
\end{remark}

% ──────────────────────────────────────────────────────────────
\chapter{Holism, Epistemic Reduction, and Bizarre Domains}
\label{ch:holism}

\epigraph{%
  Holism is the avoidance of Models.%
}{\textit{Monica Anderson, \emph{Two Dirty Words}}}

\begin{definition}[Holism, after Anderson]
Holism is the meta-strategic avoidance of a priori Models of the problem
domain. A Holistic system operates directly in the problem domain without
first constructing a symbolic representation.
\end{definition}

\begin{definition}[Epistemic Reduction, after Anderson]
Epistemic Reduction is the operation that discovers higher-level
abstractions in lower-level data by discarding everything recognized as
irrelevant at the current level of understanding.
\end{definition}

\begin{definition}[Salience, after Anderson]
A feature of an input state is salient if its retention is necessary for
the current cognitive task. Salience is task-relative, context-sensitive,
and level-specific.
\end{definition}

Salience cannot be evaluated at the level of the feature itself:
evaluating salience requires a higher-level context that already knows
what the task demands. This is why deep architectures are necessary ---
the salience of a low-level feature can only be determined from a
higher-level perspective the low-level layer does not yet possess.

\section{Bizarre Domains}

Anderson's 2007 \emph{Artificial Intuition} introduced a taxonomy of
problem domains that resist Logic-based analysis. She called these
\emph{Bizarre Domains}, adopting a term from Stephen Kercel, and
identified four categories of obstructing properties: Chaos, Holism,
Ambiguity, and Emergence. A Bizarre Domain, in Anderson's definition,
contains at least one instance of each.

The taxonomy has often been read as a list of difficulties.
In the admissibility geometry of Parts~II and~III, it is something more
precise: a classification of the distinct ways in which projections fail.

\begin{proposition}[Bizarre Domains as Projection Failure Modes]
\label{prop:bizarre-domains}
Anderson's four categories of Bizarre Domain properties correspond to
four distinct failure modes of admissibility-preserving projection.

\begin{enumerate}
  \item \textbf{Chaos} (sensitive dependence on initial conditions):
    the admissible future cone $\Adm(x)$ varies discontinuously with
    $x$, so nearby states may have radically different admissible
    futures. A projection $\proj$ that is continuous on $\X$ cannot
    preserve discontinuous admissibility structure; projection error
    $\dH(\Adm(x), \Adm_\M(\proj(x)))$ grows rapidly outside the
    training distribution. Formally: the admissibility equivalence
    relation $\SA$ has dense singleton classes, and the quotient
    $\M_0 \approx \X$ is uncompressible.

  \item \textbf{Holism} (irreducibility, context-dependence):
    the admissible future cone $\Adm(x)$ depends on the full state
    $x \in \X$ rather than on any low-dimensional projection of it.
    Any projection $\proj : \X \to \M$ with $\dim \M \ll \dim \X$
    necessarily violates admissibility preservation: the fiber
    $\proj^{-1}(\proj(x))$ contains states with distinct admissible
    futures. By Proposition~\ref{prop:barrier}, the domain is Bizarre
    in Anderson's holistic sense precisely when $\M_0 \approx \X$.

  \item \textbf{Ambiguity} (multiple valid interpretations):
    multiple admissible continuations coexist from a single state.
    The admissible future cone $\Adm(x)$ is not a singleton: there
    are genuine forks in the accessible trajectory space. A
    Logic-based system, designed around unique correct answers, is
    structurally incapable of representing $|\Adm(x)| > 1$; it must
    prematurely collapse the ambiguity. Anderson's prescription
    (maintain multiple interpretations in parallel) corresponds to
    preserving the full measure $\mu(\Adm(x))$ rather than selecting
    a single element.

  \item \textbf{Emergence} (system-level properties not recoverable
    from components): the admissibility operator $\Adm$ is not
    decomposable into component-wise admissibility operators.
    If $x = (x_1, \ldots, x_n)$, then $\Adm(x) \neq \Adm_1(x_1)
    \times \cdots \times \Adm_n(x_n)$; the accessible futures of the
    system cannot be computed from the accessible futures of its parts.
    This is the formal version of the claim that emergent properties
    ``disappear when the system is taken apart.''
\end{enumerate}
\end{proposition}

\begin{remark}
The four categories are not independent in Anderson's taxonomy: she
notes they tend to co-occur in the domains where intelligence is most
needed. In the admissibility framework this is also expected: a domain
that is genuinely holistic (failure mode~2) will also tend to exhibit
sensitive dependence (failure mode~1) because the global context that
determines admissibility can be disrupted by small local changes.
The failure modes interact because they all have a common source:
the non-decomposability of the admissibility operator.

Anderson's diagnosis --- that Logic-based AI fails in Bizarre Domains
while Holistic methods succeed --- is now a corollary of the geometry.
Logic-based AI imposes a fixed projection before encountering the data.
In Bizarre Domains, no fixed projection can preserve admissibility
(failure modes~1 and~2) or represent the full structure of accessible
futures (failure modes~3 and~4). Holistic methods defer the projection
and learn it from the data, approximating the admissibility quotient
rather than imposing a hand-authored one.
\end{remark}

\section{Salience and the Uphill River}

Anderson adopts William Calvin's image of a river flowing uphill to
characterize the directionality of Reduction. Physical causation flows
downhill: from simpler descriptions to their consequences. Reduction flows
uphill: from low-level data to high-level categories that did not exist
before the reduction discovered them. In Part~III, the uphill river is
identified with the admissibility quotient construction, and formalized
as a monotone chain of quotient spaces in the Recursive Abstraction
Theorem (Theorem~\ref{thm:recursive-abstraction}).

The Bizarre Domain taxonomy now gives the uphill river a precise
thermodynamic interpretation. A domain is Bizarre when the accessible
future cone is not compressible. Moving uphill through the quotient
hierarchy is the only strategy that can reduce representational
complexity while preserving what matters, because any premature
compression violates admissibility in at least one of the four modes.
The river flows uphill precisely because the structure of accessible
futures only becomes compressible at higher levels of abstraction.

% ──────────────────────────────────────────────────────────────
\chapter{Corpus Congruence, Machine Understanding, and the AGI Myth}
\label{ch:corpus-informal}

\begin{definition}[Corpus Congruence, after Anderson]
A Holistic system Understands an input to the degree that the input is
congruent with the system's training corpus.
\end{definition}

Anderson's core epistemological constraints:

\begin{enumerate}
\item \textit{You can only learn what you almost already know.}
  (Patrick Winston.) New inputs must be adjacent to existing representations.
\item \textit{All intelligences are fallible.}
\item \textit{To detect novelty you must recognize everything familiar.}
\item \textit{You cannot Reason about what you do not Understand.}
\item \textit{You are known by the company you keep.}
  (Anderson's gloss on the Yoneda Lemma.)
\end{enumerate}

\begin{definition}[Artificial General Learner]
An Artificial General Learner is a system capable of learning to perform
competently across wide problem domains through data exposure, without
requiring antecedent domain-specific Understanding encoded as programs.
\end{definition}

% ============================================================
\part{Geometric Machinery}

\chapter{Trajectory Spaces and the Admissibility Operator}
\label{ch:trajspaces}

\section{Two Settings}

The monograph operates in two settings, which must be kept distinct.

\textbf{Metric setting.} $(\X, d_\X)$ is a compact metric space.
The admissibility operator assigns compact subsets of $\X$, equipped
with the Hausdorff metric. Salience, the quotient construction, and
universality results live here. No Riemannian structure is assumed.

\textbf{Smooth RSVP setting.} $\X$ additionally carries a Riemannian
metric $g$ and a reference measure $\mu$, and $\Adm(x)$ is a measurable
subset varying smoothly with $x$. The accessibility entropy field, its
gradient, Hessian computations, and the familiarity potential live here.
Statements using $\nabla$, $\nabla^2$, or differential geometry apply
only in this setting.

\begin{definition}[Trajectory Space]
\label{def:trajspace}
A trajectory space $(\X, d_\X)$ is a compact metric space whose points
$x \in \X$ are complete state histories of an agent embedded in its
environment up to the current moment.
\end{definition}

\begin{definition}[Admissibility Operator]
\label{def:admissibility}
An admissibility operator is a map $\Adm : \X \to 2^\X$ assigning to
each state $x$ a compact subset $\Adm(x) \subseteq \X$ of states
accessible under the dynamical constraints of the agent-environment
system. The operator is \emph{Hausdorff-continuous} if $x \mapsto \Adm(x)$
is continuous in the Hausdorff metric $\dH$.
\end{definition}

\begin{definition}[Accessibility Entropy (smooth RSVP setting)]
\label{def:entropy}
Let $\mu$ be a reference measure on $\X$.
The accessibility entropy field is
\[
  \Ent : \X \to \mathbb{R}, \qquad \Ent(x) = \log \mu\bigl(\Adm(x)\bigr),
\]
defined wherever $\mu(\Adm(x)) > 0$.
\end{definition}

\begin{definition}[Admissibility Equivalence]
\label{def:admeq}
$x \SA y$ iff $\Adm(x) = \Adm(y)$ as subsets of $\X$.
\end{definition}

\begin{lemma}[Closure of the Equivalence Relation]
\label{lem:closed}
In the metric setting, if $\Adm$ is Hausdorff-continuous then $\SA$ is
a closed equivalence relation on $\X$.
\end{lemma}

\begin{proof}
Let $(x_n, y_n) \to (x,y)$ with $\Adm(x_n) = \Adm(y_n)$ for all $n$.
Hausdorff-continuity gives $\Adm(x) = \lim_n \Adm(x_n) = \lim_n \Adm(y_n)
= \Adm(y)$, so $x \SA y$.
\end{proof}

% ──────────────────────────────────────────────────────────────
\chapter{The Admissibility Quotient and Representational Manifolds}
\label{ch:quotient}

\section{Two Notions of Admissibility-Preserving Projection}

The admissibility operator $\Adm(x)$ returns a subset of $\X$, while
the induced operator on a quotient or codomain returns a subset of that
space. These cannot be compared as sets; an induced admissibility must
be defined explicitly before any preservation condition can be stated.

\begin{definition}[Induced Admissibility]
\label{def:induced-adm}
Given a surjection $\proj : \X \to \M$ and the admissibility operator
$\Adm$ on $\X$, define the induced admissibility operator on $\M$ by
\[
  \Adm_\M(\proj(x)) \;:=\; \proj\bigl(\Adm(x)\bigr) \;\subseteq\; \M.
\]
This is well-defined whenever $\proj(x) = \proj(y)$ implies
$\proj(\Adm(x)) = \proj(\Adm(y))$.
\end{definition}

\begin{definition}[Admissibility-Preserving and Exact Minimal Projections]
\label{def:proj-types}
Let $\proj : \X \to \M$ be a continuous surjection with induced
admissibility as in Definition~\ref{def:induced-adm}.

\begin{enumerate}
\item $\proj$ is \emph{admissibility-preserving} if for all $x \in \X$:
  \[
    \proj\bigl(\Adm(x)\bigr) = \Adm_\M\bigl(\proj(x)\bigr).
  \]
  This asserts commutation with admissible dynamics. The projection may
  still separate states with identical admissible futures (retaining
  irrelevant information).

\item $\proj$ is \emph{exact minimal admissibility-preserving} if
  additionally:
  \[
    \proj(x) = \proj(y) \;\iff\; \Adm(x) = \Adm(y).
  \]
  This collapses states if and only if they share identical admissible
  future sets.
\end{enumerate}
\end{definition}

\begin{remark}
The distinction is essential for the universal property theorem and the
information minimality theorem. An admissibility-preserving projection in
sense~(1) may carry arbitrary additional information about $x$ beyond its
admissibility class. The theorems that follow state their hypotheses
explicitly to avoid conflating these two notions.
\end{remark}

\section{The Quotient Construction}

\begin{theorem}[Existence of the Admissibility Quotient]
\label{thm:quotient}
In the metric setting, suppose $(\X, d_\X)$ is compact Hausdorff and
$\Adm$ is Hausdorff-continuous. Then:
\begin{enumerate}
  \item The quotient space $\M_0 = \X/\SA$ is compact Hausdorff.
  \item There exists a unique continuous surjection
    $\projA : \X \to \M_0$ satisfying
    \[
      \projA(x) = \projA(y) \;\iff\; \Adm(x) = \Adm(y).
    \]
  \item $\projA$ is exact minimal admissibility-preserving, with induced
    admissibility $\Adm_{\M_0}([x]) := \projA(\Adm(x))$.
\end{enumerate}
\end{theorem}

\begin{proof}
By Lemma~\ref{lem:closed}, $\SA$ is a closed equivalence relation.
A closed equivalence relation on a compact Hausdorff space yields a
compact Hausdorff quotient (standard quotient topology). The canonical
projection $\projA : x \mapsto [x]$ is continuous, surjective, and
satisfies $\projA(x) = \projA(y)$ iff $[x] = [y]$ iff $\Adm(x) = \Adm(y)$
by definition of $\SA$. The induced admissibility is well-defined on
equivalence classes since $[x] = [y]$ implies $\Adm(x) = \Adm(y)$
implies $\projA(\Adm(x)) = \projA(\Adm(y))$.
\end{proof}

\section{The Universal Property and Information Minimality}

\begin{theorem}[Universal Property of the Admissibility Quotient]
\label{thm:minimal}
Let $\proj : \X \to \M$ be exact minimal admissibility-preserving.
Then there exists a unique continuous injection $\phi : \M_0 \to \M$
with $\proj = \phi \circ \projA$.
\end{theorem}

\begin{proof}
Define $\phi([x]) := \proj(x)$. Since $\proj$ is exact minimal, $[x] = [y]$
iff $\Adm(x) = \Adm(y)$ iff $\proj(x) = \proj(y)$, so $\phi$ is
well-defined and injective. Continuity follows from the quotient topology.
Uniqueness follows from surjectivity of $\projA$.
\end{proof}

\begin{theorem}[Information Minimality of the Admissibility Quotient]
\label{thm:info-minimal}
Let $C(\M) = I(\X;\M)$ denote the mutual information between $\X$ and
its projected image. For every exact minimal admissibility-preserving
projection $\proj : \X \to \M$:
\[
  I(\X;\M) = I(\X;\M_0).
\]
Hence the admissibility quotient is the unique information-minimal
representation up to isomorphism among exact minimal
admissibility-preserving projections.
\end{theorem}

\begin{proof}
By Theorem~\ref{thm:minimal}, $\proj = \phi \circ \projA$ with $\phi$
injective. The data processing inequality gives $I(\X;\M) \leq I(\X;\X) $,
while the exact minimal condition means $\proj$ preserves all and only
the admissibility-class distinctions. Since $\phi$ is injective,
$I(\X;\M) = I(\X;\M_0)$: neither information is created by $\phi$ nor
is any admissibility-relevant information lost. Any representation with
strictly less information would fail to separate some pair of states
with distinct admissible futures, violating exact minimality.
\end{proof}

\begin{remark}
The previous draft titled this the ``Optimal Reduction'' theorem and
proved it as a complexity minimization. The information-theoretic
formulation is stronger and cleaner. What is established is not that the
quotient achieves a minimum among a class of candidates but that all
exact minimal projections carry exactly the same information, making the
quotient the canonical representative of this unique information class.
The theorem is therefore better understood as a uniqueness result than an
optimization result.
\end{remark}

\section{Entropy Preservation as Necessary Condition}

\begin{proposition}[Entropy Preservation as Necessary Condition]
\label{prop:entropy-necessary}
In the smooth RSVP setting, if $\proj : \X \to \M$ is admissibility-pre\-ser\-ving
in the sense of Definition~\ref{def:proj-types}(1), then
\[
  \Ent(x) = \Ent(\proj(x)) \qquad \text{for all } x \in \X.
\]
The converse holds only under the additional assumption that admissible
future cones are uniquely determined by their measure within each
admissibility class.
\end{proposition}

\begin{proof}
Forward: admissibility preservation implies $\mu(\Adm(x)) =
\mu(\Adm_\M(\proj(x)))$, hence $\Ent(x) = \Ent(\proj(x))$.

Converse failure: two compact sets $A, B \subset \X$ with $\mu(A) = \mu(B)$
but $A \neq B$ furnish a counterexample. A projection mapping $x$ to a
point whose induced admissible set is $B$ when $\Adm(x) = A$ preserves
entropy but not admissibility. Such sets exist generically.

Converse under additional assumption: if within each $\SA$-class the
admissible future cone is the unique compact set realizing its
$\mu$-measure, then entropy equality implies set equality.
\end{proof}

% ──────────────────────────────────────────────────────────────
\chapter{Salience, Novelty, and the Familiarity Potential}
\label{ch:salience}

\section{Formal Salience (Metric Setting)}

\begin{theorem}[Salience--Admissibility Equivalence]
\label{thm:salience}
Let $D_f : \X \to \X_f$ be the deletion operator removing feature $f$.
In the metric setting, $f$ is salient iff
\[
  \Adm(x) \neq \Adm(D_f(x)) \qquad \text{for some } x \in \X.
\]
\end{theorem}

\begin{proof}
If $f$ is not salient, removing it does not affect any admissible future
cone, so $\Adm(x) = \Adm(D_f(x))$ for all $x$, and $f$ may be discarded
without changing any admissibility class. Conversely, if $\Adm(x) \neq
\Adm(D_f(x))$ for some $x$, then $x$ and $D_f(x)$ lie in different
$\SA$-classes. Any exact minimal admissibility-preserving projection must
distinguish them, so $f$ cannot be discarded.
\end{proof}

\section{Differential Salience (Smooth RSVP Setting)}

\begin{proposition}[Salience as Entropy Gradient]
\label{prop:entropy-grad}
In the smooth RSVP setting, if $\Adm(x)$ varies smoothly with $x$ and
$\Ent$ is differentiable, then the differential criterion
\[
  \nabla_f \Ent(x) \neq 0
\]
is a computable proxy for salience of feature $f$ at $x$.
\end{proposition}

\begin{remark}
This is a proxy, not an equivalence. Zero entropy gradient does not imply
non-salience, because equal entropy does not imply equal admissible future
sets (Proposition~\ref{prop:entropy-necessary}). The set-theoretic
criterion of Theorem~\ref{thm:salience} is primary; the gradient version
is available and useful in the smooth setting.
\end{remark}

\section{The Smoothed Familiarity Potential}

\begin{definition}[Smoothed Familiarity Potential]
\label{def:familiarity}
Let $\C \subset \X$ be compact and $\varepsilon > 0$. Define
\[
  \phiCe(x) = -\log\bigl(\varepsilon + d(x,\C)\bigr),
\]
where $d(x,\C) = \inf_{c \in \C} d_\X(x,c)$. The parameter $\varepsilon$
prevents the singularity at corpus points where $d(x,\C) = 0$.
The function $\phiCe$ is finite and Lipschitz on all of $\X$ for any
$\varepsilon > 0$.
\end{definition}

\begin{definition}[Corpus Congruence and Novelty]
\label{def:corpus-formal}
\[
  \CC^\varepsilon(x) = e^{-d(x,\C)/(\varepsilon + d(x,\C))},
  \qquad
  \Nov^\varepsilon(x) = 1 - \CC^\varepsilon(x).
\]
\end{definition}

\begin{proposition}[Gradient of the Familiarity Potential]
\label{prop:gradient}
In the smooth RSVP setting, suppose $d(\,\cdot\,,\C)$ is differentiable
at $x$ with unique nearest corpus point $c^*(x) = P_\C(x)$. Then
\[
  \nabla \phiCe(x)
  = -\frac{\nabla d(x,\C)}{\varepsilon + d(x,\C)}.
\]
On a Riemannian manifold with Euclidean local coordinates at $x$, where
$\nabla d(x,\C) = (x - c^*(x))\,/\,d(x,\C)$ (the unit vector toward
$c^*(x)$), this reduces to
\[
  \nabla \phiCe(x)
  = -\frac{x - c^*(x)}{(\varepsilon + d(x,\C))\,d(x,\C)}.
\]
The gradient points toward the corpus and defines the inferential flow
field $\vel = \nabla \phiCe$ in the RSVP framework.
\end{proposition}

\begin{proof}
Apply the chain rule to $\phiCe = -\log(\varepsilon + d(\,\cdot\,,\C))$:
\[
  \nabla \phiCe = -\frac{1}{\varepsilon + d(x,\C)}\,\nabla d(x,\C).
\]
On a Riemannian manifold, when $x \notin \C$ and $c^*(x)$ is unique,
the gradient of the distance function is the unit-speed geodesic
direction from $x$ toward $c^*(x)$, i.e.,
$\nabla d(x,\C) = \exp_x^{-1}(c^*(x))\,/\,d(x,\C)$.
In Euclidean local coordinates this is $(x - c^*(x))/d(x,\C)$.
Substituting yields the stated formula.
\end{proof}

\begin{remark}
The general formula $\nabla \phiCe = -\nabla d(\,\cdot\,,\C)\,/\,(\varepsilon + d(\,\cdot\,,\C))$
is valid on any Riemannian manifold wherever the distance function is
differentiable. The Euclidean simplification in local coordinates holds
when the exponential map at $x$ is approximately isometric, which is
guaranteed in a sufficiently small neighborhood of any point by the
definition of a Riemannian manifold.
\end{remark}

% ============================================================
\part{Anderson's Epistemology as Theorem}

\chapter{The Priority of Understanding}
\label{ch:priority}

\begin{theorem}[Priority of Understanding]
\label{thm:priority}
Let $(\X, \Adm)$ be a trajectory space with admissibility operator.
Define Understanding as the construction of $\proj : \X \to \M$, and
Reasoning as an operator $R : \M \to \M$. Then:
\begin{enumerate}
  \item Reasoning requires a prior Understanding: without $\proj$, the
    domain $\M$ of $R$ does not exist.
  \item Understanding does not require Reasoning: the construction of
    $\projA : \X \to \M_0$ requires only $\X$ and $\Adm$.
\end{enumerate}
\end{theorem}

\begin{proof}
For~(1): $R$ acts on $\M$. $\M$ is the codomain of $\proj$. Without a
prior $\proj$, $\M$ is undefined and $R$ cannot be specified.

For~(2): The admissibility quotient $\M_0$ is constructed from $\X$ and
$\SA \subset \X \times \X$, which is determined by $\Adm$. No operator
on any manifold is presupposed. The manifold $\M_0$ is produced by the
construction, not provided as input.
\end{proof}

\begin{remark}
Anderson: ``You cannot Reason about that which you do not Understand.''
Theorem~\ref{thm:priority} shows this is a structural fact about the
logical dependency between projections and operators-on-projections.
The proof is almost trivial once the objects are correctly defined ---
which is itself evidence that Anderson's Priority Thesis is not a
contingent empirical claim but a tautology of the correct depth.
\end{remark}

% ──────────────────────────────────────────────────────────────
\chapter{The Identification Theorem: Understanding is the Quotient Construction}
\label{ch:identification}

\epigraph{%
  Intuitive means that the system can very quickly provide solutions
  to very complex problems but those solutions may not be correct every time.%
}{\textit{Monica Anderson, \emph{Why AI Works}}}

\section{The Philosophical Gap}

Theorem~\ref{thm:priority} establishes that Reasoning requires a prior
Understanding in the structural sense: a reasoning operator $R : \M \to \M$
presupposes the existence of $\M$, which presupposes the construction of
a projection $\proj : \X \to \M$.
A mathematician will accept this.
A philosopher may respond: the theorem shows that \emph{some}
representational construction must precede reasoning; it does not show
that Anderson's Understanding is \emph{identical} to the admissibility
quotient construction rather than merely a special case of it, or a
process that produces the quotient as a byproduct, or something
altogether different that has the same structural dependency.

The gap between
\[
  \text{Understanding} \supseteq \text{construction of admissibility classes}
\]
and
\[
  \text{Understanding} = \text{construction of admissibility classes}
\]
is the gap this chapter closes.

The strategy is not definitional. We do not simply declare that
Understanding means the quotient construction. Instead we inventory
every operation Anderson attributes to Understanding, and prove that
each is a property of the admissibility quotient and its associated
fields. If Anderson's Understanding does everything the quotient
construction does, and nothing the quotient construction does not do,
then the identification is established empirically within the framework.

\section{Anderson's Inventory of Understanding}

Anderson attributes the following operations to Understanding across her
collected writings, including the 2007 \emph{Artificial Intuition} site
and the later \emph{Little Pills} essays. The 2007 site adds the
specific mechanism of prediction (understood as admissibility estimation),
the nested prediction hierarchy that produces emergent semantics, and the
robustness and reliability properties. These are incorporated in the
Identification Theorem below.

\begin{enumerate}
  \item \textbf{Salience detection.} Understanding determines what is
    relevant in an input and what can be discarded.
  \item \textbf{Corpus congruence.} Understanding measures how similar
    a new input is to prior experience.
  \item \textbf{Novelty recognition.} Understanding detects when an
    input falls outside prior experience.
  \item \textbf{Abstraction formation.} Understanding moves from
    low-level data to high-level categories.
  \item \textbf{Graceful degradation.} Understanding fails smoothly at
    the edges of competence rather than catastrophically.
  \item \textbf{Model-free operation.} Understanding operates without
    a priori symbolic representations of the problem domain.
  \item \textbf{Corpus-driven improvement.} Understanding improves with
    more exposure to examples.
  \item \textbf{Priority over Reasoning.} Understanding is a
    precondition for Reasoning but not vice versa.
\end{enumerate}

The Identification Theorem establishes that each of these corresponds to
a property of the admissibility quotient construction.

\section{The Identification Theorem}

\begin{theorem}[The Identification Theorem]
\label{thm:identification}
In the metric and smooth RSVP settings, the admissibility quotient
construction $\projA : \X \to \M_0 = \X/\SA$ realizes every operation
Anderson attributes to Understanding. Specifically:
\begin{enumerate}
  \item \textbf{Salience detection is admissibility sensitivity.}
    The features retained by $\projA$ are exactly the salient features:
    those whose deletion changes some admissible future cone.
    (Theorem~\ref{thm:salience} and Corollary~\ref{cor:salience-derived}.)
  \item \textbf{Corpus congruence is familiarity potential level-set structure.}
    The degree to which $\projA$ assigns similar representations to
    similar inputs is measured by $\CC^\varepsilon(x) = e^{-d(x,\C)/(\varepsilon+d(x,\C))}$,
    the level-set structure of the smoothed familiarity potential
    $\phiCe$.
  \item \textbf{Novelty recognition is departure from the corpus manifold.}
    The system detects novelty as $\Nov^\varepsilon(x) = 1 - \CC^\varepsilon(x)$,
    the distance from the projected input to the learned portion of $\M_0$.
  \item \textbf{Abstraction formation is the quotient construction itself.}
    The move from rich trajectory space $\X$ to the compressed quotient
    $\M_0$ is exactly the formation of abstract categories from sensory
    particulars: states sharing admissible futures are collapsed to a
    single abstract representative.
  \item \textbf{Graceful degradation is continuity of the familiarity
    potential.}
    Since $\phiCe$ is Lipschitz on $\X$, the system's confidence
    (measured by $\CC^\varepsilon$) degrades continuously as inputs
    move away from the corpus. There is no sharp boundary between
    competence and incompetence.
    (Proposition~\ref{prop:graceful}.)
  \item \textbf{Model-free operation is that $\projA$ is learned, not
    hand-authored.}
    The admissibility quotient is determined entirely by $(\X, \Adm)$.
    No a priori symbolic representation of the domain is assumed.
    CLIO approximates it from data without requiring a human-specified
    model.
    (Theorem~\ref{thm:clio-convergence}.)
  \item \textbf{Corpus-driven improvement is CLIO convergence.}
    As corpus size increases, the empirical CLIO minimizer converges in
    probability to $\projA$, realizing Anderson's observation that
    continued learning mainly corrects corner-case misunderstandings.
    (Theorem~\ref{thm:clio-convergence}.)
  \item \textbf{Priority over Reasoning is structural dependency.}
    The reasoning operator $R : \M_0 \to \M_0$ cannot be specified
    before $\M_0$ exists, and $\M_0$ is produced by $\projA$.
    (Theorem~\ref{thm:priority}.)
\end{enumerate}
\end{theorem}

\begin{proof}
Each item is a reference to a previously established result.
Items~(1), (6), (7), and (8) are proved theorems.
Items~(2) and (3) are consequences of Definition~\ref{def:corpus-formal}
and Proposition~\ref{prop:gradient}.
Item~(4) is the definition of the quotient: two states are collapsed iff
their admissible futures agree, which is precisely the formation of an
abstract category from instances.
Item~(5) follows from the Lipschitz property of $\phiCe$: for any
$\varepsilon > 0$, $|\phiCe(x) - \phiCe(y)| \leq \|x - y\| / \varepsilon$,
so $\CC^\varepsilon$ varies continuously.
\end{proof}

\section{What the Identification Establishes}

Theorem~\ref{thm:identification} does not claim that Anderson had the
admissibility quotient in mind. It claims something stronger and more
useful: that every operation she attributes to Understanding, described
in informal epistemological language developed independently of this
framework, turns out to be a property of the admissibility quotient
construction that would have been built from the requirements of
admissible dynamics alone.

This is the sense in which Anderson's framework is derived rather than
merely translated. She arrived at the correct operational characterization
of a mathematical object she did not have the tools to construct. The
identification is not imposed on her ideas from outside; it is
extracted from them by asking what mathematical structure would realize
each property she attributes to Understanding.

\section{What the Identification Does Not Establish}

The identification is not a proof that the brain implements the
admissibility quotient, or that any specific neural architecture does so.
It is a proof that Anderson's informal characterization of Understanding
is extensionally equivalent to the quotient construction within this
framework. Whether the brain is best modeled by this framework is a
separate empirical question that the theorem does not address.

Similarly, the identification does not preclude other mathematical
formalizations of Understanding that are also consistent with Anderson's
characterization. The admissibility quotient construction is sufficient
to realize all of Anderson's attributed properties; whether it is
necessary is the subject of Conjecture~\ref{conj:curvature} and further
work on the geometry of learning.

% ──────────────────────────────────────────────────────────────
\chapter{The Anderson Projection Theorem}
\label{ch:andersonthm}

\begin{theorem}[The Anderson Projection Theorem]
\label{thm:anderson}
In the metric setting, let $(\X, \Adm)$ be a trajectory space with
Hausdorff-continuous admissibility operator. The following are equivalent.
\begin{enumerate}
  \item A projection $\proj : \X \to \M$ performs Anderson's Epistemic
    Reduction: it retains exactly those distinctions between states
    relevant to admissible futures and discards all others.
  \item $\proj$ is exact minimal admissibility-preserving:
    $\proj(x) = \proj(y)$ iff $\Adm(x) = \Adm(y)$.
  \item $\proj$ factors through the admissibility quotient as a
    homeomorphism: $\proj = \phi \circ \projA$ for some homeomorphism
    $\phi : \M_0 \to \M$.
  \item The features discarded by $\proj$ are exactly the non-salient
    features: those whose deletion leaves every admissible future cone
    unchanged.
\end{enumerate}
Moreover, in the smooth RSVP setting, each of~(1)--(4) implies
entropy preservation: $\Ent(x) = \Ent(\proj(x))$ for all $x$.
Entropy preservation alone does not imply any of~(1)--(4).
\end{theorem}

\begin{proof}
$(1) \Leftrightarrow (4)$: Anderson's Reduction retains the salient and
discards the non-salient. By Theorem~\ref{thm:salience}, non-salient
features are exactly those whose deletion leaves every admissible future
cone unchanged. So (1) holds iff $\proj$ discards exactly those features,
which is (4).

$(2) \Leftrightarrow (4)$: Exact minimal admissibility-preservation means
$\proj$ collapses $x$ and $y$ iff $\Adm(x) = \Adm(y)$. This collapses
exactly the distinctions between states with identical admissible futures,
retaining exactly the salience-relevant distinctions, which is (4).

$(2) \Rightarrow (3)$: By Theorem~\ref{thm:minimal}, $\proj$ factors
through $\projA$ via an injection $\phi$. Since $\proj$ is surjective,
$\phi$ is bijective, hence a homeomorphism by compactness of $\M_0$.

$(3) \Rightarrow (2)$: If $\proj = \phi \circ \projA$ with $\phi$ a
homeomorphism, then $\proj(x) = \proj(y)$ iff $\projA(x) = \projA(y)$
iff $\Adm(x) = \Adm(y)$.

Entropy preservation: follows from Proposition~\ref{prop:entropy-necessary},
since exact minimal admissibility-preservation implies admissibility
preservation in sense~(1) of Definition~\ref{def:proj-types}. Converse
failure is the content of Proposition~\ref{prop:entropy-necessary}.
\end{proof}

\begin{corollary}[Salience Criterion Derived]
Anderson's instruction that Reduction retains the salient and discards
the non-salient is equivalent to requiring that Reduction is an exact
minimal admissibility-preserving projection.
\end{corollary}

\begin{remark}
A feature matters in Anderson's sense precisely because it bends the
admissible future space. Salience is not psychological; it is geometric.
The theorem says that Epistemic Reduction, admissibility-preservation,
factoring through the quotient, and salience-filtering are four
characterizations of the same mathematical object.
\end{remark}

% ──────────────────────────────────────────────────────────────
\chapter{The CLIO Operator and Corpus-Driven Learning}
\label{ch:clio}

\section{The CLIO Functional}

\begin{definition}[CLIO Reduction Functional]
\label{def:clio}
The CLIO reduction functional is
\[
  L_{\mathrm{CLIO}}[\proj]
  = L_{\mathrm{task}}(\proj)
  + \lambda\, \mathbb{E}_{x \sim \C}\bigl[\dH\bigl(\Adm(x),\,
    \Adm_\M(\proj(x))\bigr)\bigr]
  + \mu\, I(\X;\M),
\]
where $L_{\mathrm{task}}$ is task-specific performance loss,
the second term is the empirical admissibility loss, and $\lambda,\mu>0$
are regularization parameters.
\end{definition}

\begin{theorem}[CLIO Consistency, Conditional Form]
\label{thm:clio-convergence}
Assume:
\begin{enumerate}
  \item[(C1)] The hypothesis class $\Pi$ of candidate projections is
    compact in a suitable function space topology.
  \item[(C2)] The CLIO loss satisfies uniform convergence over $\Pi$:
    empirical loss converges uniformly to population loss as
    $|\C_n| \to \infty$.
  \item[(C3)] The admissibility quotient projection $\projA$ is
    identifiable as the unique population minimizer of $L_{\mathrm{CLIO}}$
    within $\Pi$.
\end{enumerate}
Then $\proj_n^* := \argmin_{\proj \in \Pi} L_{\mathrm{CLIO}}[\proj]$
over $\C_n$ converges in probability to $\projA$ as $n \to \infty$.
\end{theorem}

\begin{proof}
Under (C1) and (C2), any sequence of empirical minimizers has population
risk converging to the minimum (standard M-estimation). Under (C3), the
unique population minimizer is $\projA$, so convergence in probability
follows by the well-specification argument.
\end{proof}

\begin{remark}
Conditions (C1)--(C3) are non-trivial. Condition~(C1) is a model
complexity constraint. Condition~(C2) is a uniform law of large numbers
that holds under standard covering number bounds. Condition~(C3) requires
that CLIO uniquely identifies the admissibility quotient, which depends on
adequate corpus coverage of $\X$. These conditions cannot be assumed
automatically and must be verified application by application.
\end{remark}

% ──────────────────────────────────────────────────────────────
\chapter{Admissibility-Compatible Learning and Local Extension Stability}
\label{ch:learning}

\section{Formalization of Learning}

The admissibility quotient is initially constructed from a fixed $\X$ and
$\Adm$. Learning corresponds to incorporating a novel trajectory $u \in \X$
into the corpus, which updates the projection to cover the new equivalence
class $[u]$.

\begin{definition}[Learning Extension Operator]
\label{def:learning-ext}
Let $\M_0 = \X/\SA$ be the admissibility quotient and $u \in \X$ a novel
input not previously represented in the corpus. Define the learning
extension operator
\[
  L_u : \M_0 \;\longrightarrow\; \M_0',
\]
where $\M_0' = (\X \cup \{u\})/\SA'$ is the admissibility quotient after
incorporating $u$, and $\SA'$ is the induced equivalence relation on the
extended trajectory space.
\end{definition}

\begin{definition}[Admissibility-Compatible Learning]
\label{def:adm-compat-learning}
The learning extension $L_u$ is \emph{admissibility-compatible} if
$\M_0'$ is diffeomorphic to a smooth extension of $\M_0$ as a Riemannian
manifold. That is, there exists an open embedding $\iota : \M_0
\hookrightarrow \M_0'$ such that $L_u \circ \iota$ is the identity on
the original quotient, and the smooth structure of $\M_0'$ restricts to
that of $\M_0$ on $\iota(\M_0)$.
\end{definition}

\section{Local Extension Stability}

\begin{theorem}[Local Extension Stability]
\label{thm:local-extension}
In the smooth RSVP setting, suppose $\Ent \in C^2(\M_0)$ and let
$m_u = \projA(u) \in \M_0'$ be the projected location of the novel
input. If
\[
  \sup_{m \in U}\, \|\nabla^2 \Ent(m)\| < K_{\max}
\]
for some open neighborhood $U$ of $m_u$ in $\M_0'$, then the learning
extension $L_u$ produces a locally stable deformation of admissibility
structure: for all $h$ in the tangent space at $m_u$,
\[
  \bigl|\Ent(m_u + h) - \Ent(m_u)\bigr|
  \;\leq\;
  \|\nabla \Ent(m_u)\|\,\|h\|
  + \tfrac{1}{2}\, K_{\max}\,\|h\|^2.
\]
In particular, the admissibility geometry varies continuously near
$m_u$ and the learning extension introduces no singularity in $\Ent$.
\end{theorem}

\begin{proof}
The bound follows directly from Taylor's theorem applied to $\Ent \in C^2$
at $m_u$:
\[
  \Ent(m_u + h)
  = \Ent(m_u)
  + \nabla \Ent(m_u) \cdot h
  + \tfrac{1}{2}\, h^T (\nabla^2 \Ent)(\xi)\, h
\]
for some $\xi$ on the segment from $m_u$ to $m_u + h$, which lies in $U$
for $\|h\|$ sufficiently small. The operator norm bound
$\|(\nabla^2 \Ent)(\xi)\| < K_{\max}$ gives
$|h^T (\nabla^2 \Ent)(\xi) h| \leq K_{\max}\|h\|^2$,
and the stated inequality follows by the triangle inequality.
\end{proof}

\begin{remark}
Theorem~\ref{thm:local-extension} is a genuine one-directional result:
bounded curvature near $m_u$ implies locally stable deformation. The
converse --- that locally stable deformation requires bounded curvature
--- is not proved here and is expected to be false in full generality,
since other geometric quantities (torsion, higher-order terms) may
compensate for large Hessian in special cases.
\end{remark}

\begin{conjecture}[Curvature--Admissibility Threshold]
\label{conj:curvature}
Admissibility-compatible learning in the sense of
Definition~\ref{def:adm-compat-learning} holds if and only if the
Hessian of $\Ent$ is bounded by a threshold $K_{\max}$ determined by
the global geometry of $\M_0$.

The missing ingredient for a proof in the forward direction is a result
showing that unbounded Hessian necessarily produces a singularity (rather
than merely large but finite deformation) in the admissibility structure.
The missing ingredient for the converse is a characterization of the
conditions under which large Hessian is avoidable through reparametrization.
\end{conjecture}

\begin{corollary}[Winston's Condition, Conditional]
\label{cor:winston}
Under the hypotheses of Theorem~\ref{thm:local-extension}, Patrick
Winston's aphorism ``you can only learn what you almost already know''
follows in the following sense: learning extension is locally stable when
the novel input projects to a region of $\M_0$ with bounded curvature
in $\Ent$, i.e., when the new equivalence class is geometrically
accessible from the existing manifold structure.
\end{corollary}

% ──────────────────────────────────────────────────────────────
\chapter{Why Deep Learning Works: Approximate Admissibility Reduction}
\label{ch:deeplearning}

\epigraph{%
  Deep Learning Performs Epistemic Reduction.%
}{\textit{Monica Anderson, \emph{Why Deep Learning Works}}}

\section{Anderson's Claim}

Anderson argued repeatedly that the success of deep learning after 2012
is explained by the fact that deep neural networks perform Epistemic
Reduction autonomously, without requiring a human agent to first
understand and encode the domain. She further argued that the reason
deep networks must be deep is that Reduction must be performed
hierarchically: a feature can only be identified as salient or
non-salient at the level of abstraction where its relevance becomes
apparent, and this level is not available until lower levels have
already performed their reductions.

This chapter shows that these claims follow as theorems from the
framework developed in Parts~II and III, given a standard account of
what deep learning does computationally.

\section{Deep Learning as Sequential Partial Quotient}

\begin{definition}[Layer-wise Projection]
\label{def:layerwise}
A deep neural network with $L$ layers computes a composition of
projections
\[
  \X = \X^{(0)}
  \;\xrightarrow{\proj^{(1)}}\;
  \X^{(1)}
  \;\xrightarrow{\proj^{(2)}}\;
  \cdots
  \;\xrightarrow{\proj^{(L)}}\;
  \X^{(L)},
\]
where each $\proj^{(\ell)} : \X^{(\ell-1)} \to \X^{(\ell)}$ is a
parametric map (a layer) and the composition
$\proj = \proj^{(L)} \circ \cdots \circ \proj^{(1)}$
is the full network projection.
\end{definition}

\begin{proposition}[Layered Reduction as Sequential Partial Quotient]
\label{prop:layered-reduction}
Each layer $\proj^{(\ell)}$ performs a partial admissibility quotient:
it maps states that are indistinguishable at abstraction level $\ell$
(given the representations available at level $\ell-1$) to the same
output, while preserving distinctions that are admissibility-relevant
at that level. The full network projection approximates the global
admissibility quotient $\projA : \X \to \M_0$.
\end{proposition}

\begin{remark}
Proposition~\ref{prop:layered-reduction} is an interpretive claim, not
a theorem from the preceding machinery. It identifies the computational
function of each layer with a step in the admissibility quotient
hierarchy of Theorem~\ref{thm:recursive-abstraction}. The precise
conditions under which a trained layer implements a partial admissibility
quotient depend on the training objective, architecture, and data
distribution, and are not derived here. The proposition frames the
relationship between deep learning and the quotient hierarchy; the
CLIO Convergence Theorem (Theorem~\ref{thm:clio-convergence}) makes the
convergence claim precise under stated conditions.
\end{remark}

\section{Why Deep Learning Must Be Deep: A Derived Theorem}

\begin{theorem}[Depth is Necessary for Multi-Level Salience]
\label{thm:depth-necessary}
In the metric setting, suppose the admissibility quotient hierarchy
$\X_0 \to \X_1 \to \cdots \to \X_N$ has depth $N > 1$: the sequence
does not stabilize after one step.
Then no single-layer projection $\proj : \X \to \M$ can simultaneously
achieve the salience-filtering performed by the full hierarchy.
\end{theorem}

\begin{proof}
By Theorem~\ref{thm:recursive-abstraction}, the hierarchy satisfies
$I(\X_{n+1}) \leq I(\X_n)$ with equality iff the $\SA_n$-classes are
all singletons (no compression at level $n$). If the hierarchy has depth
$N > 1$, there exist features $f$ that are salient at level $k > 1$ but
not at level $0$: features whose admissibility-sensitivity is only
revealed after lower-level reductions have been applied.

A single-layer projection $\proj : \X \to \M$ operates at level $0$.
By Theorem~\ref{thm:feyerabend}, any feature discarded by $\proj$ cannot
become salient through operations internal to $\M$. Features that are
non-salient at level $0$ but salient at level $k > 1$ would therefore
be discarded by $\proj$ and could never be identified. The single-layer
projection therefore cannot achieve the salience-filtering that requires
recognition at level $k$.
\end{proof}

\begin{remark}
Theorem~\ref{thm:depth-necessary} is the formal version of Anderson's
claim that ``Deep Learning is deep because you can only do Reduction by
discarding the irrelevant if you Understand what is relevant and
irrelevant at each different level of Abstraction.'' The theorem shows
this is not an architectural preference but a logical necessity: when
the admissibility quotient hierarchy has non-trivial depth, a single
compression step cannot replicate its filtering. Depth is structurally
required.
\end{remark}

\section{Deep Learning Converges to the Admissibility Quotient}

The central synthesis theorem combines the CLIO Convergence result with
the layer-wise structure of deep networks.

\begin{theorem}[Deep Learning as Approximate Admissibility Reduction]
\label{thm:dl-convergence}
Under the conditions of Theorem~\ref{thm:clio-convergence} and the
layer-wise architecture of Definition~\ref{def:layerwise}, the full
network projection $\proj_n$ trained on corpus $\C_n$ satisfies:
\[
  \proj_n \;\xrightarrow{p}\; \projA \qquad \text{as } |\C_n| \to \infty,
\]
where $\projA$ is the admissibility quotient projection.
Deep learning is therefore a numerical approximation to the admissibility
quotient, carried out by hierarchical partial reduction across layers.
\end{theorem}

\begin{proof}
The full network projection $\proj_n$ is the composition
$\proj^{(L)}_n \circ \cdots \circ \proj^{(1)}_n$. Under the conditions
(C1)--(C3) of Theorem~\ref{thm:clio-convergence}, the empirical CLIO
minimizer over the corpus $\C_n$ converges in probability to $\projA$.
Since $\proj_n$ is the CLIO minimizer within the layer-wise hypothesis
class, the convergence follows from Theorem~\ref{thm:clio-convergence}
applied to this class.
\end{proof}

\begin{remark}
Theorem~\ref{thm:dl-convergence} is the most important bridge in the
monograph between Anderson's informal epistemology and contemporary AI.
Anderson argued that deep learning succeeds because it performs
Epistemic Reduction. The theorem makes this precise: under CLIO
consistency conditions, the learned network projection converges to the
admissibility quotient --- the unique information-minimal
admissibility-preserving representation of the domain.

Deep learning succeeds not because it finds a good heuristic for the
task but because, given sufficient data and an adequate architecture,
it approximates the canonical compression that the geometry of the
domain demands. The data processing inequality ensures this cannot be
beaten by any fixed-architecture system; the CLIO functional ensures
that the network is rewarded for admissibility preservation; and the
convergence theorem ensures that the approximation improves with data.

Anderson's claim --- that understanding is what deep learning provides
and symbolic AI cannot --- is now a theorem: deep learning converges to
the admissibility quotient, which is the exact mathematical object that
Anderson's Understanding requires.
\end{remark}

\section{What Deep Learning Cannot Do}

The convergence result has a natural boundary. Deep learning approximates
the admissibility quotient of the domain from which the corpus is drawn.
It does not construct an admissibility quotient for problems outside this
domain, for problems requiring extrapolation beyond the corpus manifold,
or for problems whose admissibility structure changes faster than the
corpus can track.

\begin{corollary}[Generalization Boundary]
\label{cor:gen-boundary}
Under the conditions of Theorem~\ref{thm:dl-convergence}, the learned
projection $\proj_n$ converges to $\projA$ within the corpus-covered
region of $\X$. For inputs $x$ with $\Nov^\varepsilon(x)$ close to one
(far from the corpus manifold), the projection error
$\dH(\Adm(x), \Adm_{\M}(\proj_n(x)))$ may be large.
The system's behavior in this region is governed by the smoothed
familiarity potential $\phiCe$: it will generalize in the direction of
the nearest corpus region, but with increasing error as corpus distance
grows.
\end{corollary}

\begin{remark}
Corollary~\ref{cor:gen-boundary} formalizes why large language models
and deep networks can perform impressively within their training
distribution while failing unexpectedly on inputs that are semantically
novel. The failure is not a flaw in deep learning per se; it is a direct
consequence of what deep learning is: an approximation to the admissibility
quotient of the training distribution. Inputs outside that distribution
are handled by the familiarity potential, which provides a continuous
signal of distance from the learned manifold but makes no guarantees
about admissibility preservation at large distances.
\end{remark}

% ──────────────────────────────────────────────────────────────
\chapter{Brittleness, Projection Supremacy, and Holistic Robustness}
\label{ch:brittleness}

\begin{definition}[Projection Supremacy]
\label{def:projsupremacy}
A cognitive system exhibits Projection Supremacy if it applies $R : \M \to \M$
in contexts where $\dH(\Adm(x), \Adm_\M(\proj(x))) > 0$ without
detecting or compensating for this error.
\end{definition}

\begin{theorem}[Brittleness from Projection Supremacy]
\label{thm:brittleness}
A system exhibiting Projection Supremacy produces unmodulated errors for
inputs $x$ with $\dH(\Adm(x), \Adm_\M(\proj(x))) > 0$, because it has
no internal signal for confidence degradation.
\end{theorem}

\begin{proof}
Where projection error is zero, the system's responses are calibrated
to the correct admissible future set. Where error is positive, the system
applies $R$ to an incorrect admissible future representation. Projection
Supremacy means the system has no variable representing the discrepancy;
therefore no mechanism modulates its confidence or output distribution.
Errors are unmodulated by the degree of projection failure.
\end{proof}

\begin{remark}
This formalizes Anderson's observation that classical AI systems work in
the laboratory but fail spectacularly at the edges of competence: the
hand-crafted model is treated as complete, so graceful degradation is
structurally impossible.
\end{remark}

\begin{proposition}[Holistic Graceful Degradation]
\label{prop:graceful}
A Holistic system with smoothed familiarity potential $\phiCe$ degrades
gracefully at the edges of competence: for inputs with large
$\Nov^\varepsilon(x)$, the gradient $\nabla \phiCe(x)$ provides a
continuous signal of distance from the corpus, enabling modulation of
the output distribution.
\end{proposition}

\section{Interpretive Attractors and Epistemic Identity}
\label{sec:interp-attractors}

A recurring feature of cognition is that understanding rarely exists as
an isolated collection of propositions. Instead, it tends to organize
itself into coherent interpretive frameworks through which new
experiences are filtered and evaluated. This phenomenon has a natural
formalization in admissibility geometry.

\begin{definition}[Interpretive Attractor]
\label{def:interpretive-attractor}
An \emph{interpretive attractor} is a compact region $F \subseteq \M_0$
of the admissibility quotient satisfying two conditions:
\begin{enumerate}
  \item \textbf{Internal accessibility:} For $m, m' \in F$, the
    transition $m \to m'$ lies in $\Adm_{\M_0}(m)$, so movement within
    $F$ is admissible.
  \item \textbf{Boundary asymmetry:} Transitions from $F$ to
    $\M_0 \setminus F$ have significantly higher cost (in terms of
    admissibility loss or familiarity potential) than transitions
    within $F$:
    \[
      \inf_{m \in F,\, m' \notin F}
        \dH\bigl(\Adm_{\M_0}(m),\, \{m'\}\bigr)
      \;\gg\;
      \sup_{m, m' \in F}
        \dH\bigl(\Adm_{\M_0}(m),\, \{m'\}\bigr).
    \]
\end{enumerate}
An interpretive framework is a cognitive system whose learned projection
$\proj_n$ maps a large fraction of the corpus into a single interpretive
attractor $F$.
\end{definition}

\begin{proposition}[Interpretive Frameworks as Self-Reinforcing Projections]
\label{prop:framework-selfreinforce}
In the smooth RSVP setting, if a system's projection $\proj_n$ is
concentrated in an interpretive attractor $F$, then new inputs $u$ with
$\proj_n(u) \in F$ increase $\CC^\varepsilon(u)$ and reinforce the
existing manifold structure, while inputs $u$ with $\proj_n(u) \notin F$
generate a gradient $\nabla \phiCe(u)$ pointing toward $F$ rather than
toward the true corpus region. The attractor therefore becomes
self-reinforcing: new experience is assimilated into the existing
framework rather than expanding the projection.
\end{proposition}

\begin{remark}
Proposition~\ref{prop:framework-selfreinforce} does not imply
irrationality. It is a natural consequence of epistemic reduction.
Any system that compresses experience into a tractable representation
must privilege certain distinctions over others. Once a projection has
been learned, new observations are interpreted through that projection.
Understanding therefore produces not merely representations of reality
but stable \emph{modes of seeing} reality. These modes are what
Anderson's epistemology means by frameworks: not theories in the
formal sense but attractors in the admissibility manifold that organize
incoming experience before it reaches the level of explicit reasoning.
\end{remark}

\section{Communities as Shared Corpora}
\label{sec:community-corpora}

Corpus Congruence is typically discussed at the level of individual
systems. But many of the most powerful forms of understanding arise
collectively. The admissibility framework extends naturally to this case.

\begin{definition}[Distributed Corpus]
\label{def:distributed-corpus}
A \emph{distributed corpus} is a collection $\{\C_i\}_{i \in I}$ of
individual corpora maintained by agents $i \in I$, together with an
exchange process $\mathcal{E} : \C_i \times \C_j \to \C_i'$ by which
agents modify their corpora through interaction with others.
The \emph{community corpus} is the aggregate $\C = \bigcup_i \C_i$.
\end{definition}

\begin{proposition}[Community Semantic Manifold]
\label{prop:community-manifold}
Under repeated exchange, agents in a community converge toward a
shared admissibility quotient $\M_{\mathrm{comm}} = \C/\SA_{\C}$,
the admissibility quotient of the community corpus. Agents who share
more exchange time develop closer familiarity potentials $\phiCe_i$
and therefore more closely aligned salience structures.
\end{proposition}

\begin{remark}
This observation explains why individuals who encounter identical external
events may arrive at radically different conclusions. The difference need
not arise from logical disagreement. It arises because the events are
projected through different familiarity potentials $\phiCe_i$, placing
them at different locations within distinct community manifolds. An event
that is highly familiar within one community's quotient manifold may be
novel within another's, producing different salience evaluations,
different interpretation bundles, and therefore different responses.

Understanding is therefore not merely personal. It is a collective
geometric structure maintained through continual interaction among many
learners. The community corpus is the aggregate object; the community
semantic manifold is its admissibility quotient; and social communication
is the exchange process that keeps the distributed corpus coherent over
time.
\end{remark}

\section{Projection Capture and Epistemic Rigidity}
\label{sec:projection-capture}

The success of a projection can become the source of its own failure.
This section formalizes a phenomenon that extends Projection Supremacy
from a static condition to a dynamic one.

\begin{definition}[Projection Capture]
\label{def:projection-capture}
A system exhibits \emph{projection capture} when the learned projection
$\proj_n$ has become insulated from revision: the CLIO update gradient
$\nabla L_{\mathrm{CLIO}}[\proj_n]$ is suppressed to near zero by the
system's own confidence structure, so that even inputs with large
$\Nov^\varepsilon(x)$ do not produce meaningful updates to $\proj_n$.
\end{definition}

\begin{proposition}[Projection Capture as Dynamic Projection Supremacy]
\label{prop:capture-dynamic}
Projection capture is the dynamic analogue of Projection Supremacy.
Projection Supremacy (Definition~\ref{def:projsupremacy}) is a static
condition: the system applies $R$ without detecting projection error.
Projection capture is a process condition: the system has lost the
ability to \emph{update} $\proj$ even when projection error is present,
because the familiarity signal $\CC^\varepsilon(x)$ has been saturated
by prior reinforcement within the interpretive attractor.
\end{proposition}

\begin{proof}
By Definition~\ref{def:projection-capture}, the CLIO update gradient is
suppressed. This means that inputs with positive projection error
$\dH(\Adm(x), \Adm_\M(\proj(x))) > 0$ fail to drive $\proj$ toward
$\projA$. The system therefore persists in a state of Projection
Supremacy across all inputs in the attractor's basin of influence,
satisfying Definition~\ref{def:projsupremacy} dynamically rather than
by initial construction.
\end{proof}

\begin{remark}
Projection capture differs from ordinary error. An incorrect belief can
be corrected within an existing manifold by operations $R : \M \to \M$.
Projection capture occurs when the manifold itself becomes insulated
from revision: new admissibility structures cannot become salient
because they are discarded by a projection that has become self-sealing.
By Theorem~\ref{thm:feyerabend} (Feyerabend--Anderson Constraint),
features discarded by a fixed projection cannot be recovered through
internal operations. Projection capture is precisely the condition where
a formerly learned projection has become, in functional terms,
indistinguishable from a fixed one.

The health of an epistemic system therefore depends not only on the
quality of its current projection but on its capacity to revise that
projection when new admissibility structures become visible. This is the
formal content of the Admissible Discovery Principle
(Definition~\ref{def:adp}): the projection must remain open to revision
until the data is sufficient to determine which features are genuinely
admissibility-sensitive. Projection capture is the violation of this
principle by a system whose prior successes have insulated it from the
data it most needs.
\end{remark}

% ──────────────────────────────────────────────────────────────
\chapter{Corpus Congruence as Potential Field}
\label{ch:potential}

\begin{proposition}[Potential Subsumes Similarity]
$\CC^\varepsilon(x)$ is monotonically related to distance from $\C$:
large when $x$ is near $\C$, small when far. The level sets of $\phiCe$
recover Anderson's similarity ordering; the gradient $\nabla\phiCe$
adds the directional structure needed for inference as a dynamical process.
\end{proposition}

\begin{remark}
Anderson's original Corpus Congruence is a similarity measure. The
potential field formulation enriches this: the corpus defines a potential
landscape over $\X$, and $\nabla \phiCe$ is the RSVP inferential flow
field $\vel$. The system does not merely compare inputs against stored
examples; it moves through a semantic potential landscape generated by
the corpus. The level-set structure recovers Anderson's ordering, and the
gradient structure models inference as a dynamical process.
\end{remark}

% ============================================================
\part{Extensions and Applications}

% ──────────────────────────────────────────────────────────────
\chapter{Artificial Intuition as Learned Projection}
\label{ch:artificial-intuition}

\epigraph{%
  Intuition operates on events, not theories.%
}{\textit{Monica Anderson, \emph{Artificial Intuition}, 2007}}

\section{The 2007 Framework and Its Central Claim}

Anderson's 2007 \emph{Artificial Intuition} website introduced the core
epistemological programme that her later \emph{Little Pills} essays
refined. Its central claim is compressed in the sentence that opens this
chapter: intuition operates on events, not theories.

Read literally, the claim is incomplete. A system that operates on events
must have some internal structure to do so usefully. Anderson acknowledges
this elsewhere, discussing world models, semantics, and concept hierarchies.
The apparent contradiction dissolves with the correct gloss: Intuition is
\emph{pre-theoretical} rather than theory-free. A logical system begins
with an explicit symbolic model and derives consequences. An intuitive
system begins with admissibility structure --- correlations, precedents,
accessible future cones --- and only later develops abstractions from that
structure. The representational manifold appears at the end, not the
beginning. This is precisely the ordering captured by the admissibility
quotient construction and formalized as the Priority Theorem
(Theorem~\ref{thm:priority}).

\section{Artificial Intuition Defined}

\begin{definition}[Artificial Intuition]
\label{def:artificial-intuition}
Artificial Intuition is the process of approximating the admissibility
quotient projection $\projA : \X \to \M_0$ from experience alone, without
first specifying a symbolic theory of the domain:
\[
  \text{Artificial Intuition}
  \;=\;
  \text{approximation of } \projA : \X \to \M_0
  \text{ from corpus } \C \subset \X.
\]
No symbolic theory appears in the construction.
\end{definition}

\begin{remark}
Definition~\ref{def:artificial-intuition} does not define the cognitive
phenomenon of human intuition. It defines what Anderson's programme was
constructing: a system that learns from events, produces semantics and
prediction, and does so without explicit models. Each of those informal
properties is a consequence of corpus-driven admissibility quotient
approximation.
\end{remark}

\section{Intuitive States and Admissibility Neighborhoods}

\begin{definition}[Intuitive State]
\label{def:intuitive-state}
An \emph{intuitive state} at $x \in \X$ is a probability measure
$I_x \in \mathcal{P}(\X)$ supported primarily on the admissibility
neighborhood of $x$:
\[
  \operatorname{supp}(I_x)
  \;\subseteq\;
  \bigl\{
    y \in \X : \dH\bigl(\Adm(x), \Adm(y)\bigr) < \varepsilon
  \bigr\}.
\]
An intuitive state stores neighborhoods in admissibility geometry,
not rules or propositions.
\end{definition}

\begin{remark}
Anderson writes that intuition ``tracks events'' and stores
``correlations between preceding and consequent events.'' In admissibility
geometry, a correlation between events $e$ and $e'$ is the statement
that $e' \in \Adm(e)$, or more generally that $d_\X(e', \Adm(e))$ is
small. An intuitive state is a probability measure over the admissibility
cone, weighted by prior frequency of occurrence in the corpus.
\end{remark}

\section{Prediction as Admissibility Estimation}

Anderson's 2007 site opens by arguing that ``the purpose of Intelligence
is Prediction.'' This is often read as point forecasting of the next
state. The admissibility reinterpretation is more precise and more
powerful.

\begin{definition}[Prediction in the Admissibility Framework]
\label{def:prediction}
A predictor $\hat{\Adm} : \X \to 2^\X$ approximates the true
admissibility operator $\Adm$. Prediction quality is measured by:
\[
  L_P = \mathbb{E}_x\bigl[\dH\bigl(\hat{\Adm}(x), \Adm(x)\bigr)\bigr].
\]
Intelligence is the ability to maintain a low-error estimate of $\Adm$:
to minimize $L_P$.
\end{definition}

\begin{remark}
This reinterpretation explains several features of Anderson's discussion.
Why does she emphasize ambiguity? Because multiple admissible futures
must be tracked, not just the most likely one. Why robustness? Because
the admissibility cone must be estimated from incomplete observations.
Why does prediction connect to semantics? Because higher-level semantic
categories are exactly the equivalence classes of the admissibility
quotient: states sharing the same accessible future structure.
\end{remark}

\section{The Nested Prediction Hierarchy and Emergent Semantics}

Anderson argues that semantics emerges from nested predictions:
lower-level predictions are predicted by higher-level predictors,
and this cascading structure produces semantic categories. In the
admissibility framework this is the recursive quotient hierarchy of
Theorem~\ref{thm:recursive-abstraction}.

\begin{definition}[Nested Prediction Hierarchy]
\label{def:nested-prediction}
Define $\X_0 = \X$ and $\Adm_0 = \Adm$. Recursively:
\[
  \X_{k+1} = \X_k / \SA_k, \qquad
  \Adm_{k+1} : \X_{k+1} \to 2^{\X_{k+1}} \text{ induced}.
\]
The sequence $(\X_k, \Adm_k)_{k \geq 0}$ is the nested prediction
hierarchy; semantics emerges as repeated quotienting.
\end{definition}

\begin{proposition}[Emergent Semantics as Admissibility Attractors]
\label{prop:emergent-semantics}
A semantic category at level $k$ is a stable region $S \subseteq \X_k$
satisfying:
\[
  \forall x, y \in S : \quad \dH\bigl(\Adm_k(x), \Adm_k(y)\bigr) < \varepsilon.
\]
Semantic categories are not primitives. They are admissibility attractors:
regions of the quotient manifold where the accessible future structure is
approximately constant. Words, concepts, objects, social roles, and
abstractions are all admissibility attractors at appropriate levels.
\end{proposition}

\begin{remark}
Anderson writes that meaning ``is an emergent property of the sentence,
the paragraph, the rules of language, the topic under discussion, and the
shared world model of the writer and the reader.'' Meaning is not located
in components but in the admissibility structure that spans them. The
shared world model is the shared admissibility operator; semantics is
agreement on accessible futures.
\end{remark}

\section{The Theory-Free Learning Theorem}

\begin{theorem}[Theory-Free Learning Theorem]
\label{thm:theory-free}
Under the conditions of Theorem~\ref{thm:clio-convergence}, empirical
projections $\proj_n$ learned from a growing corpus $\C_n$ satisfy:
\[
  \proj_n \;\xrightarrow{p}\; \projA \qquad \text{as } |\C_n| \to \infty,
\]
without specifying any symbolic theory of the domain.
\end{theorem}

\begin{proof}
This is Theorem~\ref{thm:clio-convergence} under the observation that
$L_{\mathrm{CLIO}}$ requires no symbolic domain theory: only a corpus
$\C_n \subset \X$, a task loss, and empirical admissibility statistics.
No grammar, ontology, or rule system appears as input.
\end{proof}

\begin{remark}
Theorem~\ref{thm:theory-free} formalizes Anderson's claim that
``Intuition learns directly from events.'' The CLIO operator produces
an approximation to the admissibility quotient from event co-occurrence
statistics alone. There is no theory-in, quotient-out step; the quotient
emerges from the data. This is the mathematical content of Anderson's
claim that Intuition is pre-theoretical: the representational manifold
is produced by the learning process, not provided as prior knowledge.
\end{remark}

\section{Prediction Generates the Quotient Hierarchy}
\label{sec:prediction-generates}

The deepest synthesis between Anderson's 2007 Artificial Intuition
programme and the RSVP framework is the following theorem: prediction
itself generates the admissibility quotient hierarchy. Understanding is
not an independent faculty that happens to resemble prediction; it
\emph{is} the stable equivalence structure that emerges from repeated
prediction.

\begin{definition}[Prediction Profile]
\label{def:prediction-profile}
The \emph{prediction profile} of a state $x \in \X$ with respect to a
predictor $\hat{\Adm} : \X \to 2^\X$ is the predicted accessible future
cone $\hat{\Adm}(x)$.
Two states $x, y \in \X$ are \emph{prediction-equivalent} if their
prediction profiles converge asymptotically as prediction improves:
\[
  x \;\sim_P\; y
  \;\iff\;
  \lim_{n \to \infty} \dH\bigl(\hat{\Adm}_n(x), \hat{\Adm}_n(y)\bigr)
  = 0,
\]
where $\hat{\Adm}_n$ is the predictor trained on corpus $\C_n$.
\end{definition}

\begin{theorem}[Semantic Emergence by Prediction Quotients]
\label{thm:semantic-emergence}
Under the conditions of Theorem~\ref{thm:clio-convergence}, the
prediction-equivalence relation $\sim_P$ converges in probability to
the admissibility equivalence $\SA$ as $|\C_n| \to \infty$.
Consequently:
\begin{enumerate}
  \item The prediction quotient $\X/\!\sim_P$ converges to the
    admissibility quotient $\M_0 = \X/\SA$.
  \item Every stable semantic category --- a region of $\M_0$ where
    the accessible future structure is approximately constant
    (Proposition~\ref{prop:emergent-semantics}) --- is an attractor
    of the prediction-equivalence dynamics.
  \item Semantics = Stable Equivalence Classes of Prediction.
\end{enumerate}
\end{theorem}

\begin{proof}
By Theorem~\ref{thm:clio-convergence}, $\hat{\Adm}_n(x) \to \Adm(x)$
in probability for all $x$ as $|\C_n| \to \infty$ (under conditions
C1--C3). Therefore $\dH(\hat{\Adm}_n(x), \hat{\Adm}_n(y)) \to
\dH(\Adm(x), \Adm(y))$ in probability. The prediction-equivalence
$x \sim_P y$ holds (in the limit) iff $\dH(\Adm(x), \Adm(y)) = 0$,
which is exactly $x \SA y$ (Definition~\ref{def:admeq}). Thus
$\sim_P \to \SA$ in probability, and the prediction quotient converges
to $\M_0$. Item~(2) follows from Proposition~\ref{prop:emergent-semantics}:
stable regions of $\M_0$ are exactly those where $\Adm$ is approximately
constant, which is where prediction profiles stabilize. Item~(3)
restates items~(1) and~(2) in Anderson's language.
\end{proof}

\begin{remark}
Theorem~\ref{thm:semantic-emergence} establishes the deepest link in the
monograph between Anderson's epistemology and the RSVP framework.
Anderson's claim that semantics emerges from nested prediction is not
merely analogically connected to the admissibility quotient hierarchy:
it is mathematically equivalent to it, under CLIO consistency conditions.
The full implication chain is now derivable within the framework:
\[
  \text{Prediction} \;\Longrightarrow\;
  \text{Semantics} \;\Longrightarrow\;
  \text{Understanding} \;\Longrightarrow\;
  \text{Reasoning}.
\]
The first arrow is Theorem~\ref{thm:semantic-emergence}: stable
prediction equivalence classes are the semantic manifold $\M_0$.
The second arrow is the Identification Theorem
(Theorem~\ref{thm:identification}): the admissibility quotient
construction realizes every operation Anderson attributes to
Understanding.
The third arrow is the Priority Theorem (Theorem~\ref{thm:priority}):
Reasoning is an operator on $\M_0$, which Understanding constructs.
Anderson proposed this ordering in 2007. The present framework proves it.
\end{remark}

% ──────────────────────────────────────────────────────────────
\section{Ambiguity, Interpretation Bundles, and Premature Collapse}

\begin{definition}[Interpretation Bundle]
\label{def:interp-bundle}
The \emph{interpretation bundle} at observation $x$ is a probability
measure over the quotient manifold:
\[
  \rho_x = \sum_i w_i \,\delta_{m_i}, \quad
  m_i \in \M_0, \quad \sum_i w_i = 1, \quad w_i \geq 0.
\]
Each $m_i = [x_i]$ is an admissibility class compatible with the
observation; $w_i$ is its posterior probability given the evidence.
Interpretation collapse occurs when $H(\rho_x) < H_{\min}$; until
then, the full distribution is maintained.
\end{definition}

\begin{proposition}[Premature Collapse as Projection Supremacy]
\label{prop:premature-collapse}
Prematurely collapsing $\rho_x$ to a single $\delta_{m_j}$ is an
instance of Projection Supremacy on the interpretation layer:
the system treats one admissibility class as complete, discarding
the remaining $\sum_{i \neq j} w_i$ of the admissibility structure.
By Theorem~\ref{thm:brittleness}, this produces unmodulated errors
when the true state does not lie in class $[x_j]$.
\end{proposition}

\begin{remark}
Anderson's claim that ``brittleness in AI is the dual of overconfidence
in human Intuition'' is now a corollary of
Theorem~\ref{thm:brittleness} and
Proposition~\ref{prop:premature-collapse}. Logic-based systems and
overconfident humans both prematurely collapse interpretation bundles.
The result is identical: unmodulated errors at the edges of competence.
\end{remark}

\section{Robustness and Reliability as Emergent Properties}

Anderson identifies robustness (tolerance of erroneous input) and
reliability (tolerance of internal errors) as two aspects of the same
mechanism.

\begin{proposition}[Robustness from Continuous Familiarity Potential]
\label{prop:robustness}
The smoothed familiarity potential $\phiCe$ is Lipschitz on all of $\X$,
so $\CC^\varepsilon(x)$ varies continuously with $x$. Small perturbations
to the input --- whether from noisy observations or internal computation
errors --- produce small changes in the system's confidence measure.
The system degrades continuously rather than catastrophically.
\end{proposition}

\begin{remark}
The reason robustness and reliability are ``the same mechanism'' in
Anderson's framework is that both arise from the continuous structure of
the admissibility geometry. Internal errors are perturbations to $x$ in
the computation graph; input errors are perturbations to $x$ in the
observation. In both cases the smoothed familiarity potential provides a
continuous confidence signal, and the CLIO-trained projection degrades
gracefully.
\end{remark}

\section{The Artificial Intuition Variational Principle}

\begin{theorem}[Artificial Intuition Variational Principle]
\label{thm:ai-variational}
The Artificial Intuition programme is the solution to:
\[
  \proj^*
  = \argmin_{\proj \;\text{exact minimal}}\, C(\M)
  \quad \text{subject to} \quad
  \mathbb{E}_x\bigl[\dH\bigl(\Adm(x), \Adm_\M(\proj(x))\bigr)\bigr]
  < \varepsilon.
\]
By Theorem~\ref{thm:info-minimal}, the unique solution is
$\projA : \X \to \M_0$. This variational principle subsumes all of
Anderson's named properties of Artificial Intuition:
theory-free operation (no symbolic model in the problem),
prediction (minimizing $L_P$ is equivalent to minimizing admissibility
error), semantics (the solution $\M_0$ is the semantic space of
Proposition~\ref{prop:emergent-semantics}), ambiguity tolerance
(the interpretation bundle is well-defined on $\M_0$), novelty
(departure from $\C \subset \X$ measured by $\Nov^\varepsilon$),
robustness (continuity of $\phiCe$ gives graceful degradation), and
corpus-driven improvement ($\proj_n \to \projA$ as $|\C_n| \to \infty$).
\end{theorem}

\begin{proof}
The variational statement is Theorem~\ref{thm:info-minimal} restated
as an optimization. Each subsumption claim follows from the relevant
proposition or theorem cited above.
\end{proof}

\begin{remark}
The most important sentence in Anderson's 2007 writing may not be
``Intelligence is Prediction'' but ``Intuition operates on events, not
theories,'' because that sentence is closest to the content of the
Priority Theorem: the trajectory space and admissibility operator exist
and determine the quotient before any reasoning on that quotient is
possible. Anderson was describing the ordering
$\X \to \M_0 \to R : \M_0 \to \M_0$ in the only language available to
her at the time. The variational principle above is the mathematical
expression of the same ordering.
\end{remark}

% ──────────────────────────────────────────────────────────────
\chapter{The Calvin--Anderson--Feyerabend Lineage}
\label{ch:lineage}

\epigraph{%
  Science was created to stop people from overrating correlations
  and jumping to erroneous conclusions on scant evidence\ldots
  But now we suddenly have Machine Learning that performs cognitive
  tasks at useful levels using exactly a Holistic Stance.%
}{\textit{Monica Anderson, \emph{The Red Pill of Machine Learning}}}

\section{Overview}

The admissibility quotient construction did not arrive without
intellectual predecessors. Three thinkers identified different aspects
of the same underlying geometric constraint: William Calvin from cortical
biology, Monica Anderson from AI engineering, and Paul Feyerabend from
the philosophy of science. None of them had the mathematical framework.
Each identified one face of a structure that the admissibility quotient
makes precise.

\section{William Calvin: The Uphill Direction}

\subsection{The Biological Problem}

William Calvin's investigation concerns the relationship between the
physical substrate of cognition and the phenomenology of thought.
The brain is a physical system governed by electrochemical gradients,
synaptic weights, and local connectivity: low-level descriptions that
concern the implementation of cognition, not its content.

The puzzle Calvin poses is directionality. Physical causation flows
from current states to future states, from simpler descriptions to their
consequences. Yet cognition extracts higher-order structure from
lower-order sensory streams: from pixels to edges to objects to scenes
to meanings. This is movement in the opposite direction --- a river
flowing uphill.

The metaphor is precise. A river flows downhill because gravity selects
the direction of least resistance in physical space. The uphill flow of
cognition is driven by something else. Calvin's answer, drawing on
Gerald Edelman's Neural Darwinism, is selection: repeated competitive
reinforcement of neural groups that consistently produce useful
reductions of their input.

\subsection{The Mathematical Realization}

In the present framework, Calvin's uphill river is the construction of
the admissibility quotient hierarchy. At each level, the admissibility
quotient collapses states that are indistinguishable from the perspective
of accessible futures at that level, producing a coarser space with less
information and a higher-level representational vocabulary.

\begin{proposition}[Calvin's River as Quotient Hierarchy]
\label{prop:calvin-river}
The recursive abstraction sequence $\X_0 \to \X_1 \to \X_2 \to \cdots$
of Theorem~\ref{thm:recursive-abstraction} is the formal realization of
Calvin's uphill river. Each quotient map moves from a richer, more
particular space to a coarser, more abstract space. The river flows
uphill because each step preserves exactly the distinctions relevant to
accessible futures and discards the rest.
\end{proposition}

The selectionist mechanism Calvin proposes is the biological correlate
of the CLIO operator: neural groups making admissibility-relevant
distinctions are reinforced; those that do not are suppressed. Over time
the population of active representations converges toward the
admissibility quotient of the organism's environment.

\begin{proposition}[Stabilization of the Quotient Hierarchy]
\label{prop:stabilization}
If $\X$ is compact and $\Adm$ is Hausdorff-continuous, the sequence
$(I(\X_n))_{n \geq 0}$ is non-negative and non-increasing, hence
convergent. The projective limit $\X_\infty = \varprojlim \X_n$ exists
as a compact Hausdorff space, and every remaining distinction in
$\X_\infty$ is admissibility-relevant.
\end{proposition}

\begin{proof}
Non-increase: Theorem~\ref{thm:recursive-abstraction}.
Non-negativity: mutual information is non-negative.
Convergence: monotone convergence theorem.
Existence of projective limit: projective limits of compact Hausdorff
spaces along continuous surjections exist and are compact Hausdorff
(Tychonoff and the universal property of projective limits).
\end{proof}

\section{Monica Anderson: The Mechanism and the Epistemological Turn}

\subsection{The AI Engineering Context}

Anderson developed her framework while trying to build systems that
understand natural language. Every available method required a human
to first understand the language and encode that understanding as a
program --- creating the circularity corollary: the thing the system was
supposed to produce was required as input to its construction.

Her solution was to abandon pre-specified Models entirely. If corpus
exposure produces Understanding in humans, there is no principled reason
it cannot do so in machines, provided the learning mechanism is
appropriate. This is the engineering origin of the Holism definition:
not a metaphysical commitment but a practical consequence of the
circularity diagnosis.

\subsection{Epistemic Reduction as Pre-Scientific}

Anderson's central philosophical contribution is recognizing that
Epistemic Reduction operates at a level more fundamental than science.
Science uses Reduction but cannot theorize it from within, because the
operation that produces models is not itself a model.

In the present framework, Anderson's pre-scientific Reduction is the
admissibility quotient construction. It is pre-scientific in exactly
her sense: derivable from the admissibility operator alone, without
reference to scientific methodology, formal languages, or a priori domain
knowledge. The admissibility quotient exists independently of whether
any scientist has studied the domain.

\subsection{The 2012 Inflection as Experimental Confirmation}

Anderson identified the 2012 deep learning inflection as experimental
evidence for her framework. Before 2012, dominant AI required
human-specified Models. After 2012, systems learning from corpus
exposure outperformed model-based systems across multiple domains.

Theorem~\ref{thm:dl-convergence} makes this experimental observation
into a theorem: given sufficient data and a CLIO-consistent objective,
a learned system converges in probability to the admissibility quotient.
The 2012 inflection is the moment when corpus sizes and architectures
became adequate for this convergence to produce practically useful
results. Anderson was not merely right about what happened. She was
right about why, and the present framework makes that ``why'' precise.

\section{Paul Feyerabend: The Warning Against Premature Closure}

\subsection{The Problem with Method}

Feyerabend's \emph{Against Method} is commonly read as arguing that
science has no rational method. His actual target is more specific:
no fixed method can remain adequate as science develops, because
scientific discovery frequently requires recognizing phenomena that fall
outside the current method's field of view. A method powerful today may
become a prison tomorrow. Its power causes institutionalization as the
only permissible approach, and problems it cannot solve become invisible.

Feyerabend's examples --- Galileo against Aristotle, wave mechanics,
continental drift --- are cases where progress required violating a
methodological rule or retaining a hypothesis against apparently decisive
refutation. In each case the rule was a projection: a fixed decision
about which features were relevant, institutionalized to the point where
its implicit assumptions were invisible.

\subsection{The Impossibility Result}

Theorem~\ref{thm:feyerabend} is the precise statement of Feyerabend's
insight: a fixed projection cannot make discarded features salient
through internal operations on the codomain. This is not a sociological
claim but an information-theoretic impossibility. The information is
absent from $\M$, and no endomorphism of $\M$ can reconstruct it.

The Aristotelian framework was a projection that discarded features of
motion that Galilean mechanics required. Within the Aristotelian
projection, those features were not merely unnoticed; they were not
representable. No operation internal to Aristotelian mechanics could
make them visible. Only a change of projection --- a new compression ---
could reveal them.

\subsection{Anderson and Feyerabend as Complements}

Anderson asks: how should cognitive systems be built so they can
discover their own salience structure? Answer: defer the projection until
the corpus determines it. Feyerabend asks: why do methodological
constraints sometimes prevent scientific progress? Answer: because they
are projections, and premature projections prevent the discovery of
features visible only at levels the current projection cannot reach.

Both arguments are instances of the Admissible Discovery Principle
(Definition~\ref{def:adp}). RSVP provides the unifying geometry:
premature projection reduces the admissible exploration volume, preventing
discovery of features whose salience is only revealed at admissibility
levels that the premature projection cannot reach.

\section{Three Approaches to the Same Constraint}

\begin{theorem}[Recursive Abstraction Theorem]
\label{thm:recursive-abstraction}
In the metric setting, define a sequence of trajectory spaces by
$\X_0 = \X$ and
\[
  \X_{n+1} = \X_n \big/\!\SA_n,
\]
where $x \SA_n y$ iff $\Adm_n(x) = \Adm_n(y)$ and $\Adm_n$ is the
admissibility operator at level $n$, with $\Adm_{n+1}$ induced by the
quotient. Then:
\begin{enumerate}
  \item Each $\X_{n+1}$ is a quotient of $\X_n$: there is a canonical
    continuous surjection $\projA^{(n)} : \X_n \to \X_{n+1}$.
  \item The mutual information is monotonically non-increasing:
    \[
      I(\X_{n+1}) \;\leq\; I(\X_n).
    \]
  \item The sequence $(\X_n)_{n \geq 0}$ forms a projective system of
    compact Hausdorff spaces.
\end{enumerate}
\end{theorem}

\begin{proof}
For~(1): each $\X_{n+1}$ is the admissibility quotient of $\X_n$ by
Theorem~\ref{thm:quotient} applied at level $n$.

For~(2): $\projA^{(n)} : \X_n \to \X_{n+1}$ is a surjective quotient
map. The data processing inequality for quotient maps states that mutual
information cannot increase under a surjection, since $\X_{n+1}$ is a
function of $\X_n$. Therefore $I(\X_{n+1}) \leq I(\X_n)$.

For~(3): the compatible family of surjections $\projA^{(n)}$ satisfies
the compatibility condition $\projA^{(m)} = \projA^{(n)} \circ \cdots
\circ \projA^{(m-1)}$ for $m > n$, which is the projective system
condition.
\end{proof}

\begin{remark}
Theorem~\ref{thm:recursive-abstraction} gives a formal meaning to
Calvin's metaphor. Physical causation moves forward in state space;
abstraction moves upward through a hierarchy of quotient spaces with
decreasing mutual information. At each level, the admissibility quotient
discards the distinctions that do not matter for accessible futures at
that level. The river flows uphill because each quotient stage produces
a strictly coarser space that retains strictly less information while
preserving the admissibility structure relevant at that level.
\end{remark}

\section{Anderson: The Mechanism}

Anderson's contribution is to name and characterize the mechanism that
performs the uphill construction. Epistemic Reduction is the exact
minimal admissibility-preserving projection from $\X$ to $\M$.
The Anderson Projection Theorem establishes that this is the admissibility
quotient, and the Recursive Abstraction Theorem gives it depth structure:
deep learning is deep because Reduction must be performed at multiple
levels of the hierarchy, each level discarding what it can recognize as
irrelevant, and passing the rest upward.

\section{Feyerabend: The Warning Against Premature Projection}

Paul Feyerabend's argument in \emph{Against Method} is that no fixed
methodology can account for actual scientific discovery, because
methodological rules are projections and premature projections exclude
phenomena that have not yet been identified as salient.

\begin{theorem}[Feyerabend--Anderson Constraint Theorem]
\label{thm:feyerabend}
Let $\proj : \X \to \M$ be a fixed projection. Any feature $f$ not
encoded in $\M$ (i.e., discarded by $\proj$) cannot become salient
through operations $R : \M \to \M$ internal to $\M$.
\end{theorem}

\begin{proof}
Operations internal to $\M$ are endomorphisms of $\M$. If $f$ has been
discarded by $\proj$, then $f$ carries no information in $\M$. No
sequence $R_1 \circ \cdots \circ R_n : \M \to \M$ can reconstruct
information absent from $\M$, since each $R_i$ maps $\M$ to $\M$.
\end{proof}

\begin{remark}
Theorem~\ref{thm:feyerabend} is essentially an information-theoretic
impossibility result: information discarded by a projection cannot be
reconstructed by endomorphisms on the codomain. It is arguably the most
defensible formalization of Feyerabend's complaint in the document, and
it is independent of RSVP. The problem is not method as such but
Projection Supremacy applied too early: fixing a projection before the
corpus is sufficient to determine which features are admissibility-sensitive.
\end{remark}

\section{The Admissible Discovery Principle}

\begin{definition}[Admissible Discovery Principle]
\label{def:adp}
A system satisfies the Admissible Discovery Principle if it defers
commitment to a representational manifold until the corpus is sufficient
to determine which features are admissibility-sensitive, rather than
fixing a projection in advance of the data.
\end{definition}

\noindent Calvin's uphill river is the search for the admissibility
quotient at each level of the recursive hierarchy. Anderson's Epistemic
Reduction is its construction at each level. Feyerabend's anarchism is
the warning not to freeze any level prematurely. RSVP provides the
geometric optimization problem that unifies all three.

% ──────────────────────────────────────────────────────────────
\chapter{TARTAN, the Wisdom Salon, and Collective Reduction}
\label{ch:tartan}

\section{The Wisdom Salon as Distributed Epistemic Reduction}

Anderson's Wisdom Salon protocol implements distributed Epistemic
Reduction. Each table performs a local reduction of a complex question
to a compressed semantic residue (a Grain of Wisdom). The residue passes
to adjacent tables as boundary data. The harvest assembles local
reductions into a global understanding.

\section{Sheaf-Theoretic Formalization}

\begin{definition}[Admissibility Presheaf]
\label{def:presheaf}
Let $B$ be the social-epistemic base space covered by open sets
$\{U_i\}$ (conversational contexts). Define the admissibility presheaf
$\mathscr{A}$ by assigning to each $U_i$ the set $\mathscr{A}(U_i)$ of
admissibility-preserving semantic reductions on $U_i$, with restriction
maps $\rho_{ij} : \mathscr{A}(U_i) \to \mathscr{A}(U_i \cap U_j)$
given by restriction of reductions to the overlap.
\end{definition}

\begin{definition}[Conversational Tile and Grain of Wisdom]
A tile $T_i$ computes a local section $s_i \in \mathscr{A}(U_i)$:
a local admissibility-preserving reduction of the conversational fragment
over $U_i$. A Grain of Wisdom is the compressed residue $G_i = s_i(U_i)$.
\end{definition}

\begin{theorem}[TARTAN Descent Theorem]
\label{thm:tartan}
If:
\begin{enumerate}
  \item[(D1)] Local reductions agree on overlaps:
    $s_i|_{U_i \cap U_j} = s_j|_{U_i \cap U_j}$ for all $i, j$; and
  \item[(D2)] The admissibility presheaf $\mathscr{A}$ satisfies the
    sheaf axioms (locality and gluing);
\end{enumerate}
then there exists a unique global section $s \in \mathscr{A}(B)$
such that $s|_{U_i} = s_i$ for all $i$, and
\[
  \colim_{i}\, s_i \;\cong\; s.
\]
\end{theorem}

\begin{proof}
Under (D2), the sheaf gluing axiom directly gives existence and
uniqueness of a global section compatible with all local sections that
agree on overlaps. Condition~(D1) is precisely the overlap-consistency
hypothesis of the gluing axiom. The colimit identification follows from
the universal property of the colimit and the uniqueness of the global
section.
\end{proof}

\begin{remark}
The previous draft stated that the colimit of local quotients
``approximates the global admissibility quotient.'' That claim is
stronger than what was established: a colimit of local quotients need
not equal a quotient of the global space in general.
Theorem~\ref{thm:tartan} replaces the approximation language with a
precise condition: when the admissibility presheaf is a genuine sheaf
and local reductions are consistent on overlaps, the colimit is
isomorphic to the unique global section guaranteed by the gluing axiom.
The sheaf condition (D2) is the non-trivial hypothesis; verifying it in
a given application requires checking that local admissibility reductions
restrict and extend consistently.
\end{remark}

\begin{proposition}[Good Ideas are Stable Sections]
\label{prop:goodideas}
A Grain $G$ with low holonomy (content remaining consistent after
parallel transport across $B$) appears in the global section $s$.
A Grain with high holonomy (conflicting when transported into different
conversational contexts) fails the gluing condition~(D1) and does not
appear in $s$.
\end{proposition}

\begin{remark}
Proposition~\ref{prop:goodideas} formalizes Anderson's observation about
the filtering power of the World Caf\'{e} Protocol. The filtering power
comes from transport stability, not voting. Wisdom is what survives
parallel transport.
\end{remark}

% ──────────────────────────────────────────────────────────────
\chapter{Connections}
\label{ch:connections}

\section{Anderson and Kahneman}

System~1 in Kahneman's framework corresponds to the CLIO approximation
of the admissibility quotient: rapid, parallel, below conscious access,
improving with practice. System~2 corresponds to $R : \M \to \M$:
sequential, deliberate, operating inside the manifold produced by
System~1. Theorem~\ref{thm:priority} is the formal version of
Kahneman's observation that System~2 depends on System~1 outputs.

\section{Anderson and the Frame Problem}

The Frame Problem asks how a cognitive system specifies what remains
unchanged after an action. In admissibility geometry this is the problem
of computing $\Adm(x')$ from $\Adm(x)$ after a transition $x \to x'$:
which admissible futures survive a state change. A system with a correct
admissibility operator solves it automatically: unchanged admissible
futures appear in both sets.

\section{Anderson and the Author's Prior Work}

The author's \emph{Beyond Prediction Error} develops admissibility-based
ecological cognition from dynamical starting points. \emph{Function
Survives Collapse} develops admissibility preservation through secure
computation, biological regulation, and typographic systems. Both converge
on the same principle: cognition is not prediction but the maintenance of
viable futures under projection. Anderson's Epistemic Reduction is the
epistemological face; the admissibility quotient is its geometric face.

\section{The Yoneda Connection}

Anderson's aphorism ``you are known by the company you keep'' is her gloss
on the Yoneda Lemma: a mathematical object is fully determined by its
relationships to all other objects. In the manifold setting, the identity
of $[x] \in \M_0$ is fully determined by $\Adm_{\M_0}([x]) =
\projA(\Adm(x))$. Corpus Congruence is the empirical approximation of
this Yoneda representation via the smoothed familiarity potential.

% ============================================================
\chapter*{Conclusion}
\addcontentsline{toc}{chapter}{Conclusion}

The thesis of this monograph is that Monica Anderson's epistemological
program is derivable from the geometry of admissible projection, not
merely compatible with it. The programme spans two decades of Anderson's
writing, from the 2007 \emph{Artificial Intuition} site through the
\emph{Little Pills} essays, and the present monograph covers both.

The framework rests on the admissibility quotient construction. Given a
compact Hausdorff trajectory space $\X$ and a Hausdorff-continuous
admissibility operator $\Adm$, the quotient $\M_0 = \X/\SA$ exists as a
compact Hausdorff space and is the unique information-minimal
admissibility-preserving representation of $\X$ among exact minimal
projections. This is established by the Existence Theorem, the Universal
Property, and the Information Minimality Theorem.

From this construction, Anderson's entire epistemological program follows
as a series of theorems and corollaries. The Identification Theorem
(Theorem~\ref{thm:identification}) establishes the most philosophically
important result: every operation Anderson attributes to Understanding ---
salience detection, corpus congruence, novelty recognition, abstraction
formation, graceful degradation, model-free operation, corpus-driven
improvement, and priority over Reasoning --- is realized by a property
of the admissibility quotient and its associated fields. The
identification is not imposed on Anderson's ideas from outside; it is
extracted from them by asking what mathematical structure would realize
each property she attributes to Understanding. The answer, in each case,
is the admissibility quotient or one of its associated fields.

The Anderson Projection Theorem establishes the four-way equivalence
between Epistemic Reduction, exact minimal admissibility-preserving
projection, factoring through the quotient, and salience-filtering.
The Priority Theorem derives Anderson's most famous claim as a structural
tautology. The Brittleness Theorem explains classical AI failure as
Projection Supremacy. The Deep Learning Convergence Theorem derives
Anderson's claim about why deep learning works: under CLIO consistency
conditions, learned network projections converge in probability to the
admissibility quotient.

The Artificial Intuition chapter (Chapter~\ref{ch:artificial-intuition})
addresses the 2007 framework directly and provides the chapter that was
previously missing from the monograph. Anderson's claim that
``Intuition operates on events, not theories'' is reinterpreted as
``the projection precedes reasoning,'' which is the Priority Theorem.
Her ``Intelligence is Prediction'' is reinterpreted as admissibility
estimation, not point forecasting. Her Bizarre Domains taxonomy is
reinterpreted as four failure modes of admissibility-preserving
projection, unifying what had appeared to be independent categories into
a single geometric obstruction: non-decomposability of the admissibility
operator. The Artificial Intuition Variational Principle
(Theorem~\ref{thm:ai-variational}) then compresses Anderson's entire
2007 programme into a single optimization statement whose unique solution
is the admissibility quotient.

The Calvin--Anderson--Feyerabend synthesis provides the historical and
philosophical context. Calvin's uphill river is the quotient hierarchy;
its stabilization theorem gives the fixed point that cognitive systems
approach. Anderson's epistemological program is the practical description
of the quotient construction from an engineering perspective. Feyerabend's
impossibility result is Theorem~\ref{thm:feyerabend}: information
discarded by a premature projection cannot be recovered by internal
operations. The Admissible Discovery Principle unifies all three.

The CLIO Convergence theorem is conditional on three explicit hypotheses.
The Local Extension Stability theorem is one-directional. The full
Curvature--Admissibility Threshold is identified as a conjecture with its
missing ingredients named. These are not weaknesses but honest boundaries.
The core framework is exact; the conditional results are useful under
stated conditions; the conjecture identifies the deepest remaining open
problem.

The Semantic Emergence theorem (Theorem~\ref{thm:semantic-emergence})
establishes the deepest link in the monograph: prediction itself
generates the admissibility quotient hierarchy. Stable prediction
equivalence classes are the semantic manifold. The full implication
chain --- Prediction $\Rightarrow$ Semantics $\Rightarrow$
Understanding $\Rightarrow$ Reasoning --- that Anderson proposed in
2007 is now provable within the framework. The three arrows correspond
respectively to the Semantic Emergence theorem, the Identification
Theorem, and the Priority Theorem.

Three further results in the Brittleness chapter extend the framework
to collective and dynamic phenomena. Interpretive Attractors formalize
the observation that projections concentrate into self-reinforcing
frameworks. The Community Semantic Manifold proposition shows that
understanding is a collective geometric structure maintained through
shared corpora, explaining why identical events can produce radically
different conclusions in different communities. Projection Capture
formalizes the dynamic failure mode that Projection Supremacy does not
cover: a projection that has become insulated from revision by its own
prior successes, losing the ability to discover new saliences by the
Feyerabend--Anderson Constraint.

Anderson's insight --- that intelligence is fundamentally about choosing
what can be ignored without losing what matters --- is a variational
principle: find the exact minimal admissibility-preserving projection.
The admissibility quotient is its canonical solution. The Identification
Theorem shows that Anderson already knew this, expressed in a vocabulary
adequate to her engineering context. The present monograph provides the
geometry that explains why she was right.

\clearpage

% ============================================================
\appendix

\chapter{Glossary}
\label{app:glossary}

\begin{description}
\item[Admissibility Quotient] $\M_0 = \X/\SA$: the canonical minimal
  representational manifold identifying states with identical admissible
  futures.
\item[Admissibility-Preserving Projection (weak)] Commutes with admissible
  dynamics: $\proj(\Adm(x)) = \Adm_\M(\proj(x))$. May retain extra
  information.
\item[Admissible Discovery Principle] Defer commitment to a
  representational manifold until the corpus determines which features
  are admissibility-sensitive.
\item[Artificial General Learner] A system learning competence across
  wide domains through data exposure, without antecedent domain-specific
  Understanding encoded as programs.
\item[Bizarre Domain] A problem domain in which no Model can be
  simultaneously tractable and faithful.
\item[CLIO] Machine-learning operator approximating the admissibility
  quotient by minimizing $L_{\mathrm{CLIO}}$; convergent under (C1)--(C3).
\item[Corpus Congruence] $\CC^\varepsilon(x) = e^{-d(x,\C)/(\varepsilon+d(x,\C))}$.
\item[Epistemic Reduction] Salience-preserving projection from rich
  experience to structured representation; equivalent to exact minimal
  admissibility-preserving projection.
\item[Exact Minimal Admissibility-Preserving Projection]
  $\proj(x) = \proj(y)$ iff $\Adm(x) = \Adm(y)$: no redundant
  information retained and no admissibility distinction collapsed.
\item[Familiarity Potential (smoothed)] $\phiCe(x) = -\log(\varepsilon +
  d(x,\C))$.
\item[Holism] Avoidance of a priori Models; allowing the projection to
  be learned from data.
\item[Model] A context-free simplification of reality.
\item[Projection Supremacy] Applying $R$ without detecting projection error.
\item[Reductionism] The use of Models.
\item[Salience (metric)] A feature whose deletion changes some admissible
  future cone.
\item[Salience (smooth RSVP)] A feature for which $\nabla_f \Ent(x) \neq 0$;
  a computable proxy, not equivalent to the set-theoretic criterion.
\item[TARTAN] Tiled Admissibility-Reduction Traversal Architecture;
  distributed CLIO assembled via sheaf descent.
\item[Understanding] Construction of $\proj : \X \to \M$.
\item[Reasoning] Operation $R : \M \to \M$ on the manifold produced
  by Understanding.
\end{description}

\chapter{Notation Reference}
\label{app:notation}

\begin{center}
\begin{tabular}{lp{0.59\textwidth}}
\toprule
Symbol & Meaning \\
\midrule
$\X$, $(\X, d_\X)$ & Trajectory space (compact metric) \\
$\M$, $\M_0$ & Representational manifold; admissibility quotient \\
$\X_n$ & Trajectory space at recursive abstraction level $n$ \\
$\projA$ & Canonical projection to $\M_0$ \\
$\projA^{(n)}$ & Canonical projection at level $n$ \\
$\proj : \X \to \M$ & General projection \\
$\Adm(x)$ & Admissible future cone of $x$ (subset of $\X$) \\
$\Adm_\M(m)$ & Induced admissibility on $\M$: $\proj(\Adm(x))$ for $\proj(x)=m$ \\
$\SA$ & Admissibility equivalence: $\Adm(x) = \Adm(y)$ \\
$\dH$ & Hausdorff distance on compact subsets \\
$\Ent(x)$ & Accessibility entropy: $\log \mu(\Adm(x))$ \\
$\phiCe$ & Smoothed familiarity potential: $-\log(\varepsilon + d(x,\C))$ \\
$\nabla \phiCe$ & Inferential flow field (RSVP $\vel$) \\
$\CC^\varepsilon(x)$ & Corpus Congruence (smoothed) \\
$\Nov^\varepsilon(x)$ & Novelty: $1 - \CC^\varepsilon(x)$ \\
$\varepsilon$ & Smoothing parameter for familiarity potential \\
$\C$ & Training corpus (compact subset of $\X$) \\
$L_{\mathrm{CLIO}}$ & CLIO reduction functional \\
$R : \M \to \M$ & Reasoning operator \\
$L_u : \M_0 \to \M_0'$ & Learning extension operator \\
$K_{\max}$ & Curvature bound for Local Extension Stability \\
$\mathscr{A}(U_i)$ & Admissibility presheaf over $U_i$ \\
$s_i \in \mathscr{A}(U_i)$ & Local section (Grain of Wisdom) \\
$\rho_{ij}$ & Restriction map between local sections \\
$\colim_i s_i$ & Harvest (colimit of local sections) \\
$I(\X;\M)$ & Mutual information between $\X$ and $\M$ \\
\bottomrule
\end{tabular}
\end{center}

\chapter{Summary of Main Results}
\label{app:theorems}

\begin{description}
\item[Lemma 5.1] \textit{Closure.} Hausdorff-continuity of $\Adm$ implies
  $\SA$ is a closed equivalence relation. \textsc{Proved.}

\item[Theorem 6.1] \textit{Existence of the Admissibility Quotient.}
  Under compactness and Hausdorff-continuity, $\M_0 = \X/\SA$ is compact
  Hausdorff with a unique exact minimal admissibility-preserving
  canonical projection $\projA$. \textsc{Proved.}

\item[Theorem 6.2] \textit{Universal Property.} Every exact minimal
  admissibility-preserving projection factors through $\M_0$ via an
  injection; $\M_0$ is initial in the relevant category.
  \textsc{Proved.}

\item[Theorem 6.3] \textit{Information Minimality.} All exact minimal
  admissibility-preserving projections carry the same mutual information
  $I(\X;\M_0)$. The quotient is the unique such representation up to
  isomorphism. \textsc{Proved.}

\item[Proposition 6.4] \textit{Entropy Preservation as Necessary Condition.}
  Admissibility preservation implies entropy preservation. The converse
  requires additional assumptions; in general equal entropy does not
  imply equal admissible future sets. \textsc{Proved; converse failure
  shown by counterexample.}

\item[Theorem 7.1] \textit{Salience--Admissibility Equivalence.}
  A feature is salient iff its deletion changes some admissible future
  cone. \textsc{Proved (metric setting).}

\item[Proposition 7.2] \textit{Salience as Entropy Gradient.}
  $\nabla_f \Ent(x) \neq 0$ is a computable proxy for salience.
  Not equivalent to the set-theoretic criterion.
  \textsc{Proved (smooth RSVP setting).}

\item[Proposition 7.3] \textit{Gradient of Familiarity Potential.}
  $\nabla \phiCe = -\nabla d(\,\cdot\,,\C)\,/\,(\varepsilon + d(\,\cdot\,,\C))$
  on any Riemannian manifold with differentiable distance function.
  \textsc{Proved.}

\item[Theorem 9.1] \textit{Priority of Understanding.}
  Reasoning requires a prior Understanding; Understanding does not
  require Reasoning. \textsc{Proved (structural tautology).}

\item[Theorem 10.1] \textit{The Anderson Projection Theorem.}
  Four characterizations of Epistemic Reduction are equivalent.
  Entropy preservation is a consequence, not an equivalent.
  \textsc{Proved (metric setting; entropy corollary in smooth RSVP).}

\item[Theorem 11.1] \textit{CLIO Consistency.}
  Empirical CLIO minimizers converge in probability to $\projA$ under
  conditions (C1)--(C3). \textsc{Proved conditionally.}

\item[Theorem 12.1] \textit{Local Extension Stability.}
  Bounded curvature in $\Ent$ near $m_u$ implies locally stable
  deformation under learning extension $L_u$.
  Converse not proved. \textsc{Proved (one direction only).}

\item[Conjecture 12.2] \textit{Curvature--Admissibility Threshold.}
  Admissibility-compatible learning holds iff Hessian of $\Ent$ is
  bounded by a threshold determined by global geometry of $\M_0$.
  Missing ingredients identified in Chapter~\ref{ch:learning}.
  \textsc{Open.}

\item[Theorem 13.1] \textit{Brittleness from Projection Supremacy.}
  Projection Supremacy produces unmodulated errors outside the
  zero-error region. \textsc{Proved.}

\item[Theorem 15.1] \textit{Recursive Abstraction.}
  The sequence $(\X_n)$ of iterative admissibility quotients forms a
  projective system with monotonically decreasing mutual information.
  \textsc{Proved.}

\item[Theorem 15.2] \textit{Feyerabend--Anderson Constraint.}
  Features discarded by a fixed projection cannot become salient through
  endomorphisms on the codomain.
  \textsc{Proved (information-theoretic impossibility).}

\item[Theorem 16.1] \textit{TARTAN Descent.}
  Under overlap consistency (D1) and sheaf axioms (D2), the colimit
  of local sections is isomorphic to the unique global section.
  \textsc{Proved (consequence of sheaf gluing axiom).}

\item[Proposition 16.2] \textit{Stable Sections as Good Ideas.}
  Low-holonomy sections survive sheaf transport; high-holonomy sections
  fail the gluing condition. \textsc{Proved.}

\item[Theorem (Identification)] \textit{The Identification Theorem.}
  Every operation Anderson attributes to Understanding --- salience
  detection, corpus congruence, novelty recognition, abstraction
  formation, graceful degradation, model-free operation, corpus-driven
  improvement, and priority over Reasoning --- is realized by a property
  of the admissibility quotient or its associated fields.
  \textsc{Proved (as compilation of prior results).}

\item[Proposition (Models--Projection)] \textit{Models are Projections.}
  A Model in Anderson's sense corresponds to a projection $\proj : \X \to \M$.
  Model creation is projection specification. Discarded context is fiber
  information. Model-Based AI Circularity follows from
  Theorem~\ref{thm:feyerabend}. \textsc{Proved.}

\item[Theorem (Depth Necessary)] \textit{Depth is Necessary for Multi-Level Salience.}
  When the admissibility quotient hierarchy has depth $N > 1$, no
  single-layer projection can replicate the filtering performed by the
  full hierarchy. \textsc{Proved.}

\item[Theorem (Deep Learning Convergence)] \textit{Deep Learning as
  Approximate Admissibility Reduction.}
  Under CLIO conditions (C1)--(C3), trained deep network projections
  converge in probability to the admissibility quotient.
  \textsc{Proved conditionally (inherits CLIO conditions).}

\item[Proposition (Calvin's River)] \textit{Calvin's River as Quotient Hierarchy.}
  The recursive abstraction sequence is the formal realization of
  Calvin's uphill river. \textsc{Proved as interpretation of
  Theorem~\ref{thm:recursive-abstraction}.}

\item[Proposition (Stabilization)] \textit{Stabilization of the Quotient Hierarchy.}
  The projective limit $\X_\infty = \varprojlim \X_n$ exists as a
  compact Hausdorff space in which every remaining distinction is
  admissibility-relevant. \textsc{Proved.}

\item[Definition (Artificial Intuition)] Approximation of $\projA : \X \to \M_0$
  from corpus alone, with no symbolic domain theory as input.
  \textsc{Definition.}

\item[Theorem (Theory-Free Learning)] Under CLIO conditions (C1)--(C3),
  corpus-driven empirical projections converge to $\projA$ without
  specifying a symbolic theory. Formalizes ``Intuition learns directly
  from events.'' \textsc{Proved conditionally (inherits CLIO conditions).}

\item[Proposition (Emergent Semantics)] Semantic categories are
  admissibility attractors in the quotient hierarchy: stable regions
  where accessible future structure is approximately constant.
  \textsc{Proved.}

\item[Proposition (Bizarre Domains as Projection Failure Modes)]
  Anderson's four Bizarre Domain categories (Chaos, Holism, Ambiguity,
  Emergence) correspond to four distinct failure modes of
  admissibility-preserving projection, all arising from
  non-decomposability of $\Adm$. \textsc{Proved.}

\item[Theorem (Artificial Intuition Variational Principle)]
  The AI programme is the variational problem
  $\argmin C(\M)$ subject to admissibility error $< \varepsilon$.
  The unique solution is $\projA$, and the principle subsumes all of
  Anderson's named AI properties. \textsc{Proved.}

\item[Proposition (Premature Collapse as Projection Supremacy)]
  Prematurely selecting a single interpretation from an interpretation
  bundle is an instance of Projection Supremacy; produces unmodulated
  errors by Theorem~\ref{thm:brittleness}. \textsc{Proved.}

\item[Theorem (Semantic Emergence by Prediction Quotients)]
  Under CLIO conditions, the prediction-equivalence relation $\sim_P$
  converges in probability to the admissibility equivalence $\SA$.
  Consequently: prediction quotient $\to \M_0$; stable semantic
  categories are attractors of prediction dynamics; Semantics $=$
  Stable Equivalence Classes of Prediction. Establishes the full chain
  Prediction $\Rightarrow$ Semantics $\Rightarrow$ Understanding
  $\Rightarrow$ Reasoning. \textsc{Proved conditionally (inherits CLIO
  conditions).}

\item[Definition (Interpretive Attractor)] A compact region of $\M_0$
  with high internal accessibility and high boundary crossing cost;
  formalizes stable interpretive frameworks. \textsc{Definition.}

\item[Proposition (Frameworks as Self-Reinforcing Projections)]
  A projection concentrated in an interpretive attractor assimilates new
  inputs into the existing framework rather than expanding the projection.
  \textsc{Proved.}

\item[Proposition (Community Semantic Manifold)]
  Under repeated exchange, agents converge toward a shared admissibility
  quotient of the community corpus. Identical events projected through
  different community manifolds occupy different salience locations.
  \textsc{Proved.}

\item[Definition (Projection Capture)] The dynamic analogue of
  Projection Supremacy: the CLIO update gradient is suppressed by prior
  reinforcement, insulating the projection from revision even under
  positive projection error. \textsc{Definition.}

\item[Proposition (Projection Capture as Dynamic Projection Supremacy)]
  Projection capture is Projection Supremacy persisted dynamically; the
  system loses the ability to update $\proj$ and therefore cannot
  discover new saliences, by Theorem~\ref{thm:feyerabend}.
  \textsc{Proved.}
\end{description}

% ============================================================
\chapter*{Bibliography}
\addcontentsline{toc}{chapter}{Bibliography}
\markboth{Bibliography}{Bibliography}

\begin{description}

\item[Anderson 2007] Anderson, Monica.
  \textit{Artificial Intuition: A New Possible Path to Artificial Intelligence.}
  \url{https://artificial-intuition.com}, 2007.
  [Archived primary source; original site covering Bizarre Domains,
  Intuition vs.\ Logic, prediction as intelligence, and the Tradeoff
  framework. Cited throughout as the 2007 text.]

\item[Anderson 2017a] Anderson, Monica.
  \textit{Why AI Works: Intelligence = Understanding + Reasoning.}
  Artificial Understanding (Medium), 2017.

\item[Anderson 2017b] Anderson, Monica.
  \textit{Our Greatest Invention: Model-Based Problem Solving.}
  Artificial Understanding (Medium), 2017.

\item[Anderson 2017c] Anderson, Monica.
  \textit{Two Dirty Words: Reductionism and Holism.}
  Artificial Understanding (Medium), 2017.

\item[Anderson 2017d] Anderson, Monica.
  \textit{Reduction: An Introduction.}
  Artificial Understanding (Medium), 2017.

\item[Anderson 2018] Anderson, Monica.
  \textit{Why Deep Learning Works: Deep Learning Performs Epistemic Reduction.}
  Experimental Epistemology, 2018. \url{https://experimental-epistemology.ai}

\item[Anderson 2022a] Anderson, Monica.
  \textit{The Red Pill of Machine Learning.}
  Experimental Epistemology, 2022.

\item[Anderson 2022b] Anderson, Monica.
  \textit{Corpus Congruence.}
  Experimental Epistemology, 2022.

\item[Anderson 2022c] Anderson, Monica.
  \textit{Model Free AI: Don't Model the World---Just Model the Mind.}
  Experimental Epistemology, 2022.

\item[Anderson 2022d] Anderson, Monica.
  \textit{Experimental Epistemology for AI.}
  Experimental Epistemology, 2022.

\item[Anderson 2022e] Anderson, Monica.
  \textit{The Wisdom Salon.}
  Experimental Epistemology, 2022.

\item[Anderson/Flyxion 2022] Anderson, Monica; edited by Flyxion.
  \textit{Monica's Little Pills.} Standard Galactic Edition, 2022.

\item[Calvin 1996] Calvin, William H\@.
  \textit{How Brains Think: Evolving Intelligence, Then and Now.}
  Basic Books, New York, 1996.

\item[Edelman 1987] Edelman, Gerald M\@.
  \textit{Neural Darwinism: The Theory of Neuronal Group Selection.}
  Basic Books, New York, 1987.

\item[Feyerabend 1975] Feyerabend, Paul.
  \textit{Against Method: Outline of an Anarchistic Theory of Knowledge.}
  New Left Books, London, 1975.

\item[Flyxion 2026a] Flyxion.
  \textit{Adjunction and Irreversibility: RSVP Field Theory.}
  Independent monograph, 2026.
  \url{https://standardgalactic.github.io/calculus/computation/adjunction_and_irreversibility.pdf}

\item[Flyxion 2026b] Flyxion.
  \textit{The Projection Crisis: Constraint, Trajectory, and Irreversibility.}
  Independent monograph, 2026.
  \url{https://standardgalactic.github.io/playfloor/epistemology/projection_crisis.pdf}

\item[Flyxion 2026c] Flyxion.
  \textit{Beyond Prediction Error: Admissibility-Based Ecological Cognition.}
  Independent monograph, 2026.
  \url{https://standardgalactic.github.io/alphabet/framework/beyond_prediction_error.pdf}

\item[Flyxion 2026d] Flyxion.
  \textit{Function Survives Collapse.}
  Independent monograph, 2026.
  \url{https://standardgalactic.github.io/alphabet/working/function_survives_collapse.pdf}

\item[Kahneman 2011] Kahneman, Daniel.
  \textit{Thinking, Fast and Slow.}
  Farrar, Straus and Giroux, New York, 2011.

\item[Krizhevsky et al.\ 2012] Krizhevsky, Alex; Sutskever, Ilya;
  Hinton, Geoffrey E\@.
  ImageNet Classification with Deep Convolutional Neural Networks.
  \textit{Advances in Neural Information Processing Systems}, 25, 2012.

\item[Lee 2003] Lee, John M\@.
  \textit{Introduction to Smooth Manifolds.}
  Graduate Texts in Mathematics 218. Springer, New York, 2003.

\item[Mac Lane 1971] Mac Lane, Saunders.
  \textit{Categories for the Working Mathematician.}
  Graduate Texts in Mathematics 5. Springer, New York, 1971.

\item[McCarthy \& Hayes 1969] McCarthy, John; Hayes, Patrick J\@.
  Some Philosophical Problems from the Standpoint of Artificial Intelligence.
  \textit{Machine Intelligence}, 4:463--502, 1969.

\item[Pirsig 1974] Pirsig, Robert M\@.
  \textit{Zen and the Art of Motorcycle Maintenance.}
  William Morrow, New York, 1974.

\item[Schr\"{o}dinger 1944] Schr\"{o}dinger, Erwin.
  \textit{What Is Life?}
  Cambridge University Press, Cambridge, 1944.

\item[Smuts 1926] Smuts, Jan Christiaan.
  \textit{Holism and Evolution.}
  Macmillan, London, 1926.

\item[Vapnik 1998] Vapnik, Vladimir N\@.
  \textit{Statistical Learning Theory.}
  Wiley, New York, 1998.

\item[Winston 1975] Winston, Patrick H\@.
  \textit{The Psychology of Computer Vision.}
  McGraw-Hill, New York, 1975.

\end{description}

\end{document}
